\documentclass[11pt,a4paper]{article}

% ==================== TEMEL PAKETLER (pdflatex) ====================
\usepackage[T1]{fontenc}
\usepackage[utf8]{inputenc}
\usepackage{lmodern}
\usepackage[a4paper,margin=2cm,top=2.4cm,bottom=2.4cm]{geometry}
\usepackage{xcolor}
\usepackage{colortbl}
\usepackage{array}
\usepackage{booktabs}
\usepackage{longtable}
\usepackage{tabularx}
\usepackage{listings}
\usepackage{fancyhdr}
\usepackage{amsmath}
\usepackage{amssymb}
\usepackage{hyperref}

% ==================== RENK PALETİ ====================
\definecolor{anaMavi}{HTML}{132C4A}
\definecolor{vurguMavi}{HTML}{2E6DA4}
\definecolor{acikMavi}{HTML}{EAF2F9}
\definecolor{turuncu}{HTML}{D9701E}
\definecolor{yesil}{HTML}{2E7D32}
\definecolor{acikYesil}{HTML}{E8F5E9}
\definecolor{kirmizi}{HTML}{C0392B}
\definecolor{mor}{HTML}{6C3483}
\definecolor{acikMor}{HTML}{F3E9F7}
\definecolor{acikGri}{HTML}{F4F5F7}
\definecolor{ortaGri}{HTML}{6B7280}
\definecolor{koyuGri}{HTML}{2D2D2D}
\definecolor{kodArka}{HTML}{FBFBFA}

% ==================== HYPERREF ====================
\hypersetup{
  colorlinks=true,
  linkcolor={[HTML]{2E6DA4}},
  urlcolor={[HTML]{2E6DA4}},
  pdftitle={C++ ile GLFW ve OpenGL Rehberi},
  pdfauthor={Grafik Rehberi},
  bookmarksnumbered=true
}

% ==================== BAŞLIK STİLLERİ (titlesec olmadan) ====================
\setcounter{secnumdepth}{0}
\newcounter{secc}
\newcounter{subsecc}[secc]
\newcounter{subsubsecc}[subsecc]
\newcommand{\gsec}[1]{\par\phantomsection\stepcounter{secc}%
  \setcounter{subsecc}{0}%
  \vspace{18pt}{\color{anaMavi}\Large\bfseries\thesecc.\hspace{0.5em}#1}\par
  \vspace{2pt}{\color{turuncu}\rule{\linewidth}{1.6pt}}\par\vspace{8pt}%
  \addcontentsline{toc}{section}{\protect\numberline{\thesecc}#1}\markboth{#1}{#1}}
\newcommand{\gsub}[1]{\par\phantomsection\stepcounter{subsecc}%
  \setcounter{subsubsecc}{0}%
  \vspace{11pt}{\color{vurguMavi}\large\bfseries\thesecc.\thesubsecc\hspace{0.5em}#1}\par\vspace{4pt}%
  \addcontentsline{toc}{subsection}{\protect\numberline{\thesecc.\thesubsecc}#1}}
\newcommand{\gsubsub}[1]{\par\phantomsection\stepcounter{subsubsecc}%
  \vspace{7pt}{\color{koyuGri}\normalsize\bfseries #1}\par\vspace{2pt}}

% Kısım (Part) başlığı
\newcommand{\gpart}[2]{\clearpage
  \begin{center}\vspace*{1.5cm}
  {\color{turuncu}\rule{0.7\linewidth}{1.5pt}}\\[10pt]
  {\color{ortaGri}\Large KISIM #1}\\[8pt]
  {\color{anaMavi}\Huge\bfseries #2}\\[10pt]
  {\color{turuncu}\rule{0.7\linewidth}{1.5pt}}
  \end{center}\vspace{1cm}}

% ==================== BAŞLIK / ALTBİLGİ ====================
\setlength{\headheight}{14pt}
\pagestyle{fancy}
\fancyhf{}
\renewcommand{\headrulewidth}{0.4pt}
\renewcommand{\footrulewidth}{0.4pt}
\fancyhead[L]{\small\color{ortaGri}C++ ile GLFW \& OpenGL}
\fancyhead[R]{\small\color{ortaGri}\leftmark}
\fancyfoot[C]{\small\color{ortaGri}\thepage}

% ==================== KOD LİSTELEME (C++) ====================
\lstdefinestyle{cppstyle}{
  language=C++,
  basicstyle=\ttfamily\small\color{koyuGri},
  keywordstyle=\color{vurguMavi}\bfseries,
  commentstyle=\color{yesil}\itshape,
  stringstyle=\color{turuncu},
  numberstyle=\tiny\color{ortaGri},
  numbers=left,
  numbersep=8pt,
  backgroundcolor=\color{kodArka},
  frame=single,
  framesep=6pt,
  rulecolor=\color{ortaGri!40},
  breaklines=true,
  showstringspaces=false,
  tabsize=4,
  columns=fullflexible,
  keepspaces=true,
  extendedchars=true,
  literate={ı}{{\i}}1{İ}{{\.I}}1{ş}{{\c s}}1{Ş}{{\c S}}1{ğ}{{\u g}}1{Ğ}{{\u G}}1{ç}{{\c c}}1{Ç}{{\c C}}1{ö}{{\"o}}1{Ö}{{\"O}}1{ü}{{\"u}}1{Ü}{{\"U}}1{—}{{-{}-}}2{–}{{-}}1{’}{{'}}1{“}{{"}}1{”}{{"}}1,
  morekeywords={GLFWwindow,GLFWmonitor,GLuint,GLint,GLfloat,GLchar,GLenum,
    GLsizei,GLboolean,GLbitfield,GLvoid,GLdouble,GLubyte,glm,vec2,vec3,vec4,
    mat3,mat4,quat,uint32_t,int32_t,uint8_t,size_t,nullptr,constexpr,noexcept,
    override,final,auto,uint,in,out,layout,uniform,sampler2D,std}
}
\lstset{style=cppstyle}

% ==================== KOD LİSTELEME (GLSL) ====================
\lstdefinestyle{glslstyle}{
  basicstyle=\ttfamily\small\color{koyuGri},
  keywordstyle=\color{mor}\bfseries,
  keywordstyle=[2]\color{turuncu}\bfseries,
  commentstyle=\color{yesil}\itshape,
  stringstyle=\color{turuncu},
  numberstyle=\tiny\color{ortaGri},
  numbers=left,numbersep=8pt,
  backgroundcolor=\color{acikMor!35},
  frame=single,framesep=6pt,rulecolor=\color{mor!40},
  breaklines=true,showstringspaces=false,tabsize=4,
  columns=fullflexible,keepspaces=true,extendedchars=true,
  morekeywords={void,float,int,bool,vec2,vec3,vec4,mat2,mat3,mat4,
    in,out,inout,uniform,layout,location,sampler2D,samplerCube,
    if,else,for,while,return,const,struct,discard},
  morekeywords=[2]{gl_Position,gl_FragCoord,gl_PointSize,gl_VertexID,
    texture,normalize,dot,cross,reflect,pow,max,min,mix,clamp,length},
  literate={ı}{{\i}}1{İ}{{\.I}}1{ş}{{\c s}}1{Ş}{{\c S}}1{ğ}{{\u g}}1{Ğ}{{\u G}}1{ç}{{\c c}}1{Ç}{{\c C}}1{ö}{{\"o}}1{Ö}{{\"O}}1{ü}{{\"u}}1{Ü}{{\"U}}1{—}{{-{}-}}2{–}{{-}}1{’}{{'}}1{“}{{"}}1{”}{{"}}1,
}

% Kabuk/terminal stili
\lstdefinestyle{shstyle}{
  basicstyle=\ttfamily\small\color{koyuGri},
  backgroundcolor=\color{koyuGri!6},
  frame=single,framesep=6pt,rulecolor=\color{ortaGri!40},
  breaklines=true,showstringspaces=false,tabsize=4,columns=fullflexible,
  keepspaces=true,
  literate={ı}{{\i}}1{İ}{{\.I}}1{ş}{{\c s}}1{Ş}{{\c S}}1{ğ}{{\u g}}1{Ğ}{{\u G}}1{ç}{{\c c}}1{Ç}{{\c C}}1{ö}{{\"o}}1{Ö}{{\"O}}1{ü}{{\"u}}1{Ü}{{\"U}}1{—}{{-{}-}}2{–}{{-}}1,
}

% ==================== ÖZEL KUTULAR (tcolorbox olmadan) ====================
\newcommand{\callout}[4]{\par\smallskip\noindent\begingroup
  \setlength{\fboxsep}{8pt}\setlength{\fboxrule}{0.8pt}%
  \fcolorbox{#1}{#2}{\parbox{\dimexpr\linewidth-2\fboxsep-2\fboxrule\relax}{%
    {\bfseries\color{#1}#3}\par\smallskip #4}}\endgroup\par\smallskip}
\newcommand{\bilgi}[1]{\callout{vurguMavi}{acikMavi}{Bilgi}{#1}}
\newcommand{\ipucu}[1]{\callout{yesil}{acikYesil}{İpucu}{#1}}
\newcommand{\uyari}[1]{\callout{kirmizi}{kirmizi!7}{Dikkat}{#1}}
\newcommand{\ornek}[2][Örnek]{\callout{ortaGri!90!black}{acikGri}{#1}{#2}}
\newcommand{\notk}[2][Not]{\callout{mor}{acikMor}{#1}{#2}}

% Yardımcı komutlar
\newcommand{\code}[1]{\texttt{\color{koyuGri}#1}}
\newcommand{\hd}[1]{\texttt{\color{koyuGri}\textless #1\textgreater}}
\renewcommand{\arraystretch}{1.35}

% ==================== BAŞLANGIÇ ====================
\begin{document}

% ---------- KAPAK ----------
\begin{titlepage}
\centering
\vspace*{1.4cm}
{\color{turuncu}\rule{\textwidth}{2pt}}\\[8pt]
{\color{anaMavi}\Huge\bfseries C++ ile GLFW}\\[10pt]
{\color{vurguMavi}\LARGE Bilgisayar Grafiği, Matematik ve OpenGL Rehberi}\\[8pt]
{\color{ortaGri}\large Teoriden Pencereye, Matematikten İlk Üçgene Kapsamlı Başvuru}\\[6pt]
{\color{turuncu}\rule{\textwidth}{2pt}}\\[1.4cm]

\begin{center}
\fcolorbox{vurguMavi}{acikMavi}{\parbox{0.86\textwidth}{\centering\large
Bu belge önce bilgisayar grafiğinin \textbf{nasıl çalıştığını} (grafik boru
hattı, GPU, rasterizasyon), sonra grafiğin \textbf{matematiğini} (vektörler,
matrisler, dönüşümler, projeksiyon) ve ardından \textbf{GLFW + OpenGL} ile
adım adım kod yazmayı \textbf{çok sayıda çalışan örnekle} anlatır.
Tüm kod örnekleri İngilizce isimlendirmeyle, üretim projelerindeki
gerçek stile uygun yazılmıştır.\vspace{4pt}}}
\end{center}

\vspace{0.7cm}
\begin{center}
\fcolorbox{ortaGri!50}{acikGri}{\parbox{0.86\textwidth}{\small
\textbf{Kapsam --- Kısım 1 (Grafik Temelleri):} bilgisayar grafiği nedir
\textbullet\ raster vs vektör \textbullet\ GPU \& CPU \textbullet\ grafik boru hattı
\textbullet\ koordinat uzayları \textbullet\ renk \& framebuffer \textbullet\ çift tamponlama.\vspace{3pt}

\textbf{Kapsam --- Kısım 2 (Grafik Matematiği):} vektörler \textbullet\ nokta \& çapraz çarpım
\textbullet\ matrisler \& sayısal örnekler \textbullet\ öteleme/ölçek/döndürme \textbullet\ homojen koordinatlar
\textbullet\ Model-View-Projection \textbullet\ kamera (LookAt) \textbullet\ perspektif \& ortografik projeksiyon.\vspace{3pt}

\textbf{Kapsam --- Kısım 3 (GLFW):} kurulum \& CMake \textbullet\ ilk pencere
\textbullet\ hata callback'i \textbullet\ olay döngüsü \textbullet\ GLAD \textbullet\ klavye/fare/scroll girdisi
\textbullet\ tam ekran \& monitör \textbullet\ delta time \& FPS sayacı.\vspace{3pt}

\textbf{Kapsam --- Kısım 4 (OpenGL Çizim):} GLSL derinlemesine \textbullet\ VBO/VAO/EBO
\textbullet\ tam ilk üçgen programı \textbullet\ uniform'lar \textbullet\ dokular \& çoklu doku karışımı
\textbullet\ glm ile dönüşümler \textbullet\ 3B küp \& derinlik testi \textbullet\ FPS kamera sınıfı \textbullet\ Phong aydınlatma.\vspace{3pt}

\textbf{Kapsam --- Kısım 5 (Projeler):} soyutlanmış Shader sınıfı \textbullet\ hata denetimi yardımcıları
\textbullet\ dönen renkli küp (tam program) \textbullet\ hata ayıklama tablosu \textbullet\ sonraki adımlar.\vspace{3pt}}}
\end{center}

\vfill
{\color{ortaGri}\large Hazırlanma Tarihi: 14 Temmuz 2026}\\[4pt]
{\color{ortaGri} C++17 \textbullet\ GLFW 3.x \textbullet\ OpenGL 3.3 Core \textbullet\ GLAD \textbullet\ GLM}
\vspace*{0.4cm}
\end{titlepage}

% ---------- İÇİNDEKİLER ----------
\tableofcontents
\thispagestyle{empty}
\clearpage

% #####################################################################
\gpart{1}{Bilgisayar Grafiğinin Temelleri}
% #####################################################################

% =====================================================================
\gsec{Bilgisayar Grafiği Nedir?}
% =====================================================================

Bilgisayar grafiği, sayısal verileri (noktalar, renkler, geometri) ekrandaki
piksellere dönüştürme bilimidir. Bir üçgenin köşe koordinatlarından yola çıkıp,
onu ekranda renkli bir yüzey olarak görmemizi sağlayan tüm matematiksel ve
donanımsal süreç bu başlık altındadır. GLFW ve OpenGL öğrenmeden önce
\emph{neyi} kontrol ettiğimizi anlamak, sonraki her satırı çok daha anlamlı kılar.

\gsub{Raster ve Vektör Grafik}

İki temel grafik temsili vardır. Aradaki farkı bilmek, GPU'nun neden
piksellerle çalıştığını anlamanın anahtarıdır.

\begin{center}
\begin{tabularx}{\textwidth}{@{}>{\bfseries\color{anaMavi}}l >{\raggedright\arraybackslash}X >{\raggedright\arraybackslash}X@{}}
\toprule
\rowcolor{acikMavi}\color{anaMavi}\textbf{Tür} & \color{anaMavi}\textbf{Nasıl saklanır?} & \color{anaMavi}\textbf{Örnek}\\
\midrule
Raster & Piksellerden oluşan bir ızgara; her piksel bir renk değeri tutar. Büyütünce bozulur (pikselleşme). & PNG, JPEG, ekran görüntüsü, fotoğraf.\\
\rowcolor{acikGri} Vektör & Matematiksel şekiller (doğru, eğri, çokgen). Sonsuz büyütülebilir, bozulmaz. & SVG, font, CAD çizimi.\\
\bottomrule
\end{tabularx}
\end{center}

Ekranımız bir raster aygıttır: sonunda her şey piksele dönmek zorundadır. OpenGL
ile aslında \emph{vektörel} tanımlar (köşe koordinatları) veririz; GPU bunları
\textbf{rasterizasyon} adı verilen adımda piksellere çevirir.

\gsub{Gerçek Zamanlı ve Offline Grafik}

\begin{center}
\begin{tabularx}{\textwidth}{@{}>{\bfseries\color{anaMavi}}l >{\raggedright\arraybackslash}X@{}}
\toprule
\rowcolor{acikMavi}\color{anaMavi}\textbf{Yaklaşım} & \color{anaMavi}\textbf{Açıklama}\\
\midrule
Gerçek zamanlı (real-time) & Kareyi milisaniyeler içinde üretmek gerekir (60 FPS $\approx$ 16.6 ms/kare, 144 FPS $\approx$ 6.9 ms/kare). Oyunlar, simülasyonlar. Rasterizasyon tabanlıdır ve \textbf{OpenGL/GLFW bu dünyadadır}.\\
\rowcolor{acikGri} Offline (ışın izleme) & Kare başına saniyeler/dakikalar harcanabilir. Fotogerçekçilik hedeflenir. Sinema efektleri (ray tracing / path tracing).\\
\bottomrule
\end{tabularx}
\end{center}

\bilgi{Bu rehberdeki her şey \textbf{gerçek zamanlı rasterizasyon} üzerinedir.
Amacımız kareyi mümkün olan en hızlı şekilde çizip ekrana basmaktır; bunun için
GPU'nun paralel gücünü kullanırız. 60 FPS demek, tüm sahneyi (girdi, mantık,
çizim) 16.6 milisaniyede bitirmek zorunda olmak demektir --- bu yüzden her çağrının
maliyeti önemlidir.}

\gsub{CPU ve GPU: İş Bölümü}

CPU birkaç güçlü çekirdekle sıralı işleri iyi yapar; GPU ise binlerce basit
çekirdekle \emph{aynı işlemi} milyonlarca veri üzerinde paralel yürütür. Bir
ekran 1920$\times$1080 = ~2 milyon piksel içerir; her karede hepsini boyamak
paralellik ister. İşte GPU tam bu iş için vardır.

\begin{center}
\begin{tabular}{|>{\centering\arraybackslash}p{0.44\textwidth}|>{\centering\arraybackslash}p{0.44\textwidth}|}
\hline
\rowcolor{anaMavi}\color{white}\textbf{CPU} & \color{white}\textbf{GPU}\\
\hline
\cellcolor{acikMavi} Az sayıda (4--32) güçlü çekirdek & \cellcolor{acikYesil} Binlerce basit çekirdek\\
\hline
\cellcolor{acikMavi} Karmaşık, dallı, sıralı mantık & \cellcolor{acikYesil} Basit, paralel, aynı-işlem\\
\hline
\cellcolor{acikMavi} Düşük gecikme (latency) & \cellcolor{acikYesil} Yüksek verim (throughput)\\
\hline
\end{tabular}
\end{center}

\notk[Programlama modeli]{CPU'da C++ ile \emph{host} kodu yazarız (GLFW penceresi,
oyun mantığı). GPU'da ise GLSL ile \emph{shader} yazarız (her köşe/piksel için
çalışan minik programlar). OpenGL bu ikisi arasındaki köprüdür: CPU'daki veriyi
GPU'ya yükler ve shader'ları çalıştırır. Verinin CPU'dan GPU'ya kopyalanması
pahalıdır; bu yüzden veriyi bir kez yükleyip defalarca kullanmayı hedefleriz.}

% =====================================================================
\gsec{Grafik Boru Hattı (Graphics Pipeline)}
% =====================================================================

Grafik boru hattı, köşe verisinin (vertex) ekrandaki renkli piksele dönüşene
kadar geçtiği aşamalar zinciridir. OpenGL 3.3 Core'da bu boru hattının bazı
aşamalarını biz programlarız (shader), bazıları ise sabit-fonksiyondur (GPU
donanımı yapar). Tüm OpenGL kodunuz aslında bu boru hattını beslemekten ibarettir.

\gsub{Boru Hattının Aşamaları}

\begin{center}
\begin{tabularx}{\textwidth}{@{}>{\bfseries}r >{\raggedright\arraybackslash}X c@{}}
\toprule
\rowcolor{acikMavi}\color{anaMavi}\textbf{Aşama} & \color{anaMavi}\textbf{Ne yapar?} & \color{anaMavi}\textbf{Kim?}\\
\midrule
1. Vertex verisi & VBO'dan köşe özniteliklerini (konum, renk, UV) okur. & Biz (VAO)\\
\rowcolor{acikGri} 2. Vertex Shader & Her köşeyi işler; genelde MVP ile clip uzayına taşır. & \textbf{Biz (GLSL)}\\
3. Primitive Assembly & Köşeleri üçgen/çizgi/nokta olarak gruplar. & GPU\\
\rowcolor{acikGri} 4. (Geometry Shader) & Opsiyonel: yeni geometri üretebilir. & Biz (ops.)\\
5. Rasterizasyon & Üçgeni kapsadığı piksellere (fragment) böler; öznitelikleri interpolasyonla yayar. & GPU\\
\rowcolor{acikGri} 6. Fragment Shader & Her fragment'in nihai rengini hesaplar (doku, ışık). & \textbf{Biz (GLSL)}\\
7. Testler \& Karıştırma & Derinlik testi, stencil, alpha blend uygulanır. & GPU\\
\rowcolor{acikGri} 8. Framebuffer & Sonuç piksel renk tamponuna yazılır. & GPU\\
\bottomrule
\end{tabularx}
\end{center}

\bilgi{OpenGL 3.3 Core Profile'da \textbf{en az iki shader yazmak zorunludur}:
bir \emph{vertex shader} ve bir \emph{fragment shader}. Bunlar olmadan hiçbir şey
çizilmez. Rehberin 4. kısmında bunları satır satır yazacağız.}

\gsub{Bir Üçgenin Yolculuğu}

Somut düşünelim: ekranın ortasına bir üçgen çizmek istiyoruz. Üç köşesinin
konumunu C++ dizisinde tanımlarız $\rightarrow$ VBO ile GPU belleğine yükleriz
$\rightarrow$ vertex shader her köşeyi \code{gl\_Position} olarak clip uzayına
yerleştirir $\rightarrow$ rasterizasyon üçgenin içindeki tüm pikselleri bulur
$\rightarrow$ fragment shader her pikseli (ör. turuncuya) boyar $\rightarrow$
sonuç ekrana gelir.

\ornek[Zihinsel model]{OpenGL bir "durum makinesidir" (state machine). Sen
"şu tamponu bağla", "şu shader'ı kullan", "şimdi çiz" dersin; OpenGL o anki
duruma göre boru hattını çalıştırır. Bu yüzden çağrı \emph{sırası} çok önemlidir.
Bir nesneyi \emph{bind} etmek, "sonraki komutlar bu nesneyi etkilesin" demektir.}

% =====================================================================
\gsec{Koordinat Sistemleri ve Uzaylar}
% =====================================================================

Bir 3B nesne ekrana gelene kadar birçok koordinat uzayından geçer. Bu
dönüşümlerin tamamı matris çarpımlarıdır (Kısım 2). Şimdilik hangi uzayların
olduğunu ve ne işe yaradığını kavrayalım.

\begin{center}
\begin{tabularx}{\textwidth}{@{}>{\bfseries\color{anaMavi}}l >{\raggedright\arraybackslash}X >{\raggedright\arraybackslash}X@{}}
\toprule
\rowcolor{acikMavi}\color{anaMavi}\textbf{Uzay} & \color{anaMavi}\textbf{Açıklama} & \color{anaMavi}\textbf{Dönüşüm}\\
\midrule
Local (Model) & Nesnenin kendi merkezine göre koordinatları. & ---\\
\rowcolor{acikGri} World & Nesnenin sahnedeki (dünyadaki) yeri. & Model matrisi\\
View (Kamera) & Her şeyin kameraya göre konumu. & View matrisi\\
\rowcolor{acikGri} Clip & Perspektif uygulanmış, kırpma yapılan uzay. & Projection matrisi\\
NDC & Normalized Device Coords: $[-1,1]^3$ küpü. & Perspektif bölme ($w$)\\
\rowcolor{acikGri} Screen & Piksel koordinatları (viewport). & Viewport dönüşümü\\
\bottomrule
\end{tabularx}
\end{center}

\notk[MVP zinciri]{Bir köşenin ekrandaki yeri şu çarpımla bulunur:
\[
\mathbf{v}_{clip} = \mathbf{P} \cdot \mathbf{V} \cdot \mathbf{M} \cdot \mathbf{v}_{local}
\]
Bu üç matrisi (Model, View, Projection) vertex shader'a uniform olarak
göndeririz. Kısım 2 bu matrislerin nasıl kurulduğunu, Kısım 4 ise koda nasıl
döküldüğünü anlatır.}

\gsub{Normalized Device Coordinates (NDC)}

Vertex shader'dan çıkan \code{gl\_Position} clip uzayındadır. GPU bunu $w$
bileşenine bölerek NDC'ye geçer. NDC, her ekseni $[-1, 1]$ arasında olan bir
küptür: sol alt köşe $(-1,-1)$, sağ üst $(+1,+1)$, merkez $(0,0)$. Bu küpün
dışında kalan her şey \textbf{kırpılır} (clip). En basit üçgeni doğrudan NDC'de
tanımlayarak (matris kullanmadan) çizebiliriz --- ilk örneğimizde tam olarak
bunu yapacağız.

\begin{center}
\begin{tabular}{|>{\centering\arraybackslash}p{0.3\textwidth}|>{\centering\arraybackslash}p{0.3\textwidth}|}
\hline
\rowcolor{anaMavi}\multicolumn{2}{|c|}{\color{white}\textbf{NDC Düzlemi (z'yi yok sayarsak)}}\\
\hline
\cellcolor{acikMor}$(-1,+1)$ sol üst & \cellcolor{acikMor}$(+1,+1)$ sağ üst\\
\hline
\multicolumn{2}{|c|}{\cellcolor{acikGri} merkez $(0,0)$}\\
\hline
\cellcolor{acikYesil}$(-1,-1)$ sol alt & \cellcolor{acikYesil}$(+1,-1)$ sağ alt\\
\hline
\end{tabular}
\end{center}

\uyari{OpenGL'de NDC'nin Y ekseni \textbf{yukarı} doğrudur (matematiksel gelenek),
ama ekran piksel koordinatlarında Y genelde \textbf{aşağı} doğrudur. Doku (texture)
koordinatlarında bu ters çevrim sık sık kafa karıştırır; ilerde göreceğiz.}

% =====================================================================
\gsec{Renk, Framebuffer ve Çift Tamponlama}
% =====================================================================

\gsub{Renklerin Temsili}

Ekranda renk, kırmızı-yeşil-mavi (RGB) ışığın toplamıyla oluşur (toplamsal renk
karışımı). OpenGL'de renkleri genelde $[0.0, 1.0]$ aralığında \code{float}
olarak veririz; bazen alfa (saydamlık) ile birlikte RGBA olarak.

\begin{center}
\begin{tabularx}{\textwidth}{@{}c c c c >{\raggedright\arraybackslash}X@{}}
\toprule
\rowcolor{acikMavi}\color{anaMavi}\textbf{R} & \color{anaMavi}\textbf{G} & \color{anaMavi}\textbf{B} & \color{anaMavi}\textbf{A} & \color{anaMavi}\textbf{Renk}\\
\midrule
1.0 & 0.0 & 0.0 & 1.0 & Kırmızı\\
\rowcolor{acikGri} 0.0 & 1.0 & 0.0 & 1.0 & Yeşil\\
0.0 & 0.0 & 1.0 & 1.0 & Mavi\\
\rowcolor{acikGri} 1.0 & 1.0 & 1.0 & 1.0 & Beyaz\\
0.2 & 0.3 & 0.3 & 1.0 & Koyu turkuaz (klasik arka plan)\\
\bottomrule
\end{tabularx}
\end{center}

\gsub{Framebuffer ve Çift Tamponlama}

Framebuffer, o an ekranda gösterilecek pikselleri tutan bellek bloğudur. Eğer
tek bir tampona hem çizip hem gösterirsek, yarı çizilmiş kareler görünür ve
görüntü titrer (flicker, tearing). Çözüm \textbf{çift tamponlamadır}:

\begin{center}
\begin{tabular}{|>{\centering\arraybackslash}p{0.42\textwidth}|>{\raggedright\arraybackslash}p{0.46\textwidth}|}
\hline
\rowcolor{anaMavi}\color{white}\textbf{Tampon} & \color{white}\textbf{Rol}\\
\hline
\cellcolor{acikMavi}\textbf{Back Buffer} & GPU bir sonraki kareyi buraya \emph{gizlice} çizer.\\
\hline
\cellcolor{acikYesil}\textbf{Front Buffer} & Ekranda o an görünen kare.\\
\hline
\end{tabular}
\end{center}

Kare bittiğinde iki tampon \textbf{yer değiştirir} (swap). GLFW'de bunu
\code{glfwSwapBuffers(window)} çağrısı yapar. Böylece kullanıcı hiçbir zaman
yarı çizilmiş kare görmez.

\notk[VSync]{Tamponların yer değiştirmesini ekranın yenileme hızıyla (60/144 Hz)
senkronlamak "VSync"tir. GLFW'de \code{glfwSwapInterval(1)} ile açılır; ekran
yırtılmasını (tearing) engeller ama FPS'i yenileme hızına sabitler.
\code{glfwSwapInterval(0)} ile kapatılır (FPS sınırsız, ölçüm için faydalı).}

\ipucu{Kısım 1'i bitirdin! Artık grafiğin \emph{ne} yaptığını biliyorsun:
vektörel geometriyi GPU boru hattından geçirip, koordinat uzayları arasında
dönüştürerek, çift tamponlu bir framebuffer'a rasterize etmek. Şimdi bu
dönüşümlerin \emph{matematiğine} geçiyoruz.}

% #####################################################################
\gpart{2}{Grafik Matematiği}
% #####################################################################

Bu kısım grafiğin dilini kurar. Endişelenme: her kavramı hem matematiksel
gösterimle hem de birazdan kullanacağımız \textbf{GLM} kütüphanesindeki karşılığıyla
vereceğiz. GLM (OpenGL Mathematics), GLSL sözdizimini C++'a taşıyan başlık-only
bir kütüphanedir; \code{glm::vec3}, \code{glm::mat4} gibi tipleri sağlar. Bütün
kod örneklerinde İngilizce isimler kullanacağız.

% =====================================================================
\gsec{Trigonometri: Grafiğin Açı Dili}
% =====================================================================

Trigonometri, açılar ile uzunluklar arasındaki ilişkidir ve grafiğin her yerinde
karşımıza çıkar: döndürme matrisleri, dairesel hareket, kamera yönü, dalga
efektleri, aydınlatma... Döndürmeyi gerçekten anlamak için önce sinüs ve kosinüsü
sağlam oturtmak gerekir. Bu bölüm sıfırdan kurar.

\gsub{Radyan ve Derece}

Açıyı iki birimle ölçeriz. \textbf{Derece} günlük birimdir (tam tur $=360^\circ$).
\textbf{Radyan} ise matematiğin ve tüm grafik kütüphanelerinin doğal birimidir:
bir açının radyan değeri, birim çemberde o açının gerdiği yay uzunluğudur. Tam
tur $=2\pi$ radyandır.

\[
180^\circ = \pi \text{ radyan} \qquad\Longrightarrow\qquad
\text{radyan} = \text{derece}\times\frac{\pi}{180}
\]

\begin{center}
\begin{tabularx}{\textwidth}{@{}c c c >{\raggedright\arraybackslash}X@{}}
\toprule
\rowcolor{acikMavi}\color{anaMavi}\textbf{Derece} & \color{anaMavi}\textbf{Radyan} & \color{anaMavi}\textbf{$\sin$} & \color{anaMavi}\textbf{$\cos$}\\
\midrule
$0^\circ$   & $0$        & $0$            & $1$\\
\rowcolor{acikGri} $30^\circ$  & $\pi/6$    & $0.5$          & $\approx 0.866$\\
$45^\circ$  & $\pi/4$    & $\approx 0.707$ & $\approx 0.707$\\
\rowcolor{acikGri} $60^\circ$  & $\pi/3$    & $\approx 0.866$ & $0.5$\\
$90^\circ$  & $\pi/2$    & $1$            & $0$\\
\rowcolor{acikGri} $180^\circ$ & $\pi$      & $0$            & $-1$\\
$270^\circ$ & $3\pi/2$   & $-1$           & $0$\\
\bottomrule
\end{tabularx}
\end{center}

\uyari{C++ \code{<cmath>} fonksiyonları (\code{sin}, \code{cos}, \code{tan}) ve
GLM'nin trigonometrik/döndürme fonksiyonları \textbf{radyan} bekler. Derece ile
çalışmak istersen daima \code{glm::radians(derece)} ile çevir. Bu, yeni
başlayanların bir numaralı hatasıdır --- \code{sin(90)}, sandığın $1$'i değil,
$90$ radyanın sinüsü olan $\approx 0.894$'ü verir.}

\begin{lstlisting}
#include <cmath>
#include <glm/glm.hpp>

const float PI = 3.14159265358979f;

float radians = glm::radians(90.0f);   // 1.5708 (pi/2)
float degrees = glm::degrees(PI);      // 180.0
float s = std::sin(radians);           // 1.0
float c = std::cos(radians);           // ~0.0
\end{lstlisting}

\gsub{Birim Çember: Sinüs ve Kosinüs}

En sağlam sezgi \emph{birim çemberdir}: merkezi orijinde, yarıçapı 1 olan çember.
Merkezden $\theta$ açısıyla çizilen yarıçapın çember üzerindeki ucu tam olarak
$(\cos\theta,\ \sin\theta)$ noktasıdır.

\[
\text{çember üzerindeki nokta} = (\cos\theta,\ \sin\theta)
\]

Yani: \textbf{kosinüs} noktanın X (yatay) bileşeni, \textbf{sinüs} ise Y (dikey)
bileşenidir. $\theta$ arttıkça nokta çember boyunca saat yönünün tersine döner.
Bu tek cümle, döndürme ve dairesel hareketin tamamının temelidir.

\begin{center}
\begin{tabular}{|>{\centering\arraybackslash}p{0.28\textwidth}|>{\centering\arraybackslash}p{0.28\textwidth}|>{\centering\arraybackslash}p{0.28\textwidth}|}
\hline
\rowcolor{anaMavi}\color{white}\textbf{Açı $\theta$} & \color{white}\textbf{$\cos\theta$ (X)} & \color{white}\textbf{$\sin\theta$ (Y)}\\
\hline
\cellcolor{acikMavi}$0^\circ$ (sağ) & $1$ & $0$\\
\hline
\cellcolor{acikYesil}$90^\circ$ (yukarı) & $0$ & $1$\\
\hline
\cellcolor{acikMor}$180^\circ$ (sol) & $-1$ & $0$\\
\hline
\cellcolor{turuncu!18}$270^\circ$ (aşağı) & $0$ & $-1$\\
\hline
\end{tabular}
\end{center}

\begin{lstlisting}
// Move an object along a circle of radius r (classic orbit)
float angle  = (float)glfwGetTime();   // grows over time (radians)
float radius = 2.0f;
float x = radius * std::cos(angle);    // horizontal position
float y = radius * std::sin(angle);    // vertical position
glm::vec3 orbitPos(x, y, 0.0f);        // point orbiting the origin
\end{lstlisting}

\bilgi{Dik üçgen tanımıyla köprü: bir dik üçgende $\sin\theta=\dfrac{\text{karşı}}
{\text{hipotenüs}}$, $\cos\theta=\dfrac{\text{komşu}}{\text{hipotenüs}}$,
$\tan\theta=\dfrac{\sin\theta}{\cos\theta}=\dfrac{\text{karşı}}{\text{komşu}}$.
Birim çemberde hipotenüs $=1$ olduğu için karşı kenar doğrudan $\sin$, komşu
kenar doğrudan $\cos$ olur --- iki tanım aynı şeydir.}

\gsub{Temel Özdeşlikler}

Grafik kodunda sık işine yarayacak birkaç özdeşlik:

\begin{center}
\begin{tabularx}{\textwidth}{@{}>{\raggedright\arraybackslash}X >{\raggedright\arraybackslash}X@{}}
\toprule
\rowcolor{acikMavi}\color{anaMavi}\textbf{Özdeşlik} & \color{anaMavi}\textbf{Nerede işe yarar?}\\
\midrule
$\sin^2\theta + \cos^2\theta = 1$ & Pisagor; birim vektör olduğunu doğrular, normalize kısayolu.\\
\rowcolor{acikGri} $\sin(-\theta) = -\sin\theta$ & Tek fonksiyon (simetri).\\
$\cos(-\theta) = \cos\theta$ & Çift fonksiyon (simetri).\\
\rowcolor{acikGri} $\sin(\theta) = \cos(\theta - 90^\circ)$ & Sinüs, kosinüsün kaydırılmışıdır (faz farkı).\\
Periyot $= 2\pi$ & Her ikisi de $2\pi$'de tekrarlar; dalga/animasyon döngüsü.\\
\bottomrule
\end{tabularx}
\end{center}

\ornek[$(\cos\theta,\sin\theta)$ neden birim vektördür?]{Uzunluğu hesaplayalım:
\[
\lVert(\cos\theta,\sin\theta)\rVert=\sqrt{\cos^2\theta+\sin^2\theta}=\sqrt{1}=1.
\]
İşte bu yüzden $(\cos\theta,\sin\theta)$ her zaman \emph{birim} bir yön verir ---
kamera bakış yönü üretirken (yaw/pitch) tam da bundan faydalanırız.}

\gsub{Tanjant ve atan2: Noktadan Açıya}

Bazen tersine gitmek gerekir: elimizde bir yön vektörü var, açısını istiyoruz.
Bunun için \code{atan2(y, x)} kullanılır --- sıradan \code{atan}'dan farkı, dört
bölgeyi (quadrant) doğru ayırt etmesi ve $-\pi$ ile $+\pi$ arasında doğru açıyı
vermesidir.

\begin{lstlisting}
// Direction from player to target -> angle to face it
glm::vec2 toTarget = targetPos - playerPos;
float facingAngle  = std::atan2(toTarget.y, toTarget.x); // radians, -PI..PI

// Turn it back into a facing direction vector:
glm::vec2 facing(std::cos(facingAngle), std::sin(facingAngle));
\end{lstlisting}

\uyari{Her zaman \code{atan2(y, x)} kullan, \code{atan(y/x)} değil. \code{atan}
noktanın hangi çeyrekte olduğunu bilemez (ör. $(1,1)$ ile $(-1,-1)$ aynı oranı
verir) ve $x=0$'da sıfıra bölme yaşarsın. \code{atan2} bu iki sorunu da çözer.}

\gsub{Trigonometri Neden Döndürme Matrisinin İçinde?}

Şimdi her şey birleşiyor. Bir $(x, y)$ noktasını $\theta$ açısıyla döndürünce yeni
konum şudur:
\[
x' = x\cos\theta - y\sin\theta,\qquad y' = x\sin\theta + y\cos\theta
\]
Bu iki satır tam olarak Z-ekseni döndürme matrisinin ilk iki satırıdır (bir
sonraki bölümde göreceğiz). Yani döndürme matrisindeki bütün o $\sin$/$\cos$
değerleri, "noktayı birim çember boyunca $\theta$ kadar kaydır" demenin matris
dilindeki karşılığıdır.

\begin{lstlisting}
// Rotating a 2D point by hand (what the rotation matrix does internally)
glm::vec2 rotate2D(glm::vec2 p, float theta) {
    float s = std::sin(theta), c = std::cos(theta);
    return glm::vec2(p.x * c - p.y * s,   // x' = x cos - y sin
                     p.x * s + p.y * c);  // y' = x sin + y cos
}
\end{lstlisting}

\bilgi{Trigonometri ile animasyon: \code{sin} ve \code{cos} $[-1, 1]$ arasında
yumuşakça salındığı için "nefes alma", yukarı-aşağı süzülme, yanıp sönme gibi
efektlerin doğal kaynağıdır. \code{float offset = std::sin(glfwGetTime() * hız) *
genlik;} tek satırla pürüzsüz bir salınım verir. Aynı fikir, GLSL'de dalga/su
efektlerinin de temelidir.}

\ipucu{Grafikte trigonometriyi üç yerde göreceksin: (1) \textbf{döndürme}
(matrislerdeki $\sin$/$\cos$), (2) \textbf{yön üretme} (yaw/pitch $\to$ birim
vektör, kamera bölümü), (3) \textbf{periyodik hareket} (yörünge, dalga,
salınım). Üçü de bu bölümdeki tek fikre dayanır: açı $\to$ çember üzerinde nokta.}

% =====================================================================
\gsec{Vektörler}
% =====================================================================

Vektör, hem büyüklüğü (uzunluk) hem yönü olan bir niceliktir. Grafikte konumları,
yönleri, renkleri, hızları temsil eder. 3B'de üç bileşeni vardır: $\mathbf{v} =
(x, y, z)$.

\gsub{Temel İşlemler}

\begin{center}
\begin{tabularx}{\textwidth}{@{}>{\bfseries\color{anaMavi}}l l >{\raggedright\arraybackslash}X@{}}
\toprule
\rowcolor{acikMavi}\color{anaMavi}\textbf{İşlem} & \color{anaMavi}\textbf{Formül} & \color{anaMavi}\textbf{Anlamı}\\
\midrule
Toplama & $\mathbf{a}+\mathbf{b}=(a_x{+}b_x,\dots)$ & Yer değiştirmelerin zinciri.\\
\rowcolor{acikGri} Skalerle çarpma & $s\mathbf{a}=(sa_x,\dots)$ & Uzatma/kısaltma.\\
Uzunluk & $\lVert\mathbf{a}\rVert=\sqrt{a_x^2+a_y^2+a_z^2}$ & Büyüklük (norm).\\
\rowcolor{acikGri} Normalizasyon & $\hat{\mathbf{a}}=\mathbf{a}/\lVert\mathbf{a}\rVert$ & Uzunluğu 1 olan birim yön.\\
\bottomrule
\end{tabularx}
\end{center}

\begin{lstlisting}
#include <glm/glm.hpp>
#include <glm/gtc/type_ptr.hpp>

glm::vec3 a(1.0f, 2.0f, 3.0f);
glm::vec3 b(4.0f, 5.0f, 6.0f);

glm::vec3 sum       = a + b;              // (5, 7, 9)
glm::vec3 scaled    = 2.0f * a;           // (2, 4, 6)
float     len       = glm::length(a);     // sqrt(14) ~ 3.742
glm::vec3 unit       = glm::normalize(a); // length becomes 1
\end{lstlisting}

\uyari{Bir vektörü normalize etmek onu \emph{birim vektör} yapar (uzunluğu 1).
Işık ve kamera hesaplarının neredeyse tamamı yön vektörlerinin normalize
olmasını bekler; unutursan aydınlatma yanlış çıkar. \code{glm::normalize} ile
sıfır vektörü normalize etme --- 0'a bölme NaN üretir.}

\gsub{Nokta Çarpım (Dot Product)}

İki vektörün nokta çarpımı bir \emph{skalerdir} ve aralarındaki açıyla ilgilidir:
\[
\mathbf{a}\cdot\mathbf{b}=a_xb_x+a_yb_y+a_zb_z=\lVert\mathbf{a}\rVert\,\lVert\mathbf{b}\rVert\cos\theta
\]
Her iki vektör de birim ise, nokta çarpım doğrudan $\cos\theta$'dır. Bu, grafikte
altın değerindedir: bir yüzeyin ışığa ne kadar dönük olduğunu (diffuse
aydınlatma) tam olarak bu belirler.

\begin{center}
\begin{tabularx}{\textwidth}{@{}c >{\raggedright\arraybackslash}X@{}}
\toprule
\rowcolor{acikMavi}\color{anaMavi}\textbf{$\mathbf{a}\cdot\mathbf{b}$ işareti} & \color{anaMavi}\textbf{Yorum}\\
\midrule
$>0$ & Vektörler aynı yöne bakıyor (açı $<90^\circ$).\\
\rowcolor{acikGri} $=0$ & Dik (ortogonal).\\
$<0$ & Zıt yöne bakıyor (açı $>90^\circ$).\\
\bottomrule
\end{tabularx}
\end{center}

\begin{lstlisting}
float d = glm::dot(a, b);                    // 1*4 + 2*5 + 3*6 = 32

// Angle between two directions (in radians):
glm::vec3 u = glm::normalize(a);
glm::vec3 v = glm::normalize(b);
float angle = glm::acos(glm::dot(u, v));     // 0..PI

// Is the surface facing the light? (diffuse factor)
float diffuse = glm::max(glm::dot(normal, lightDir), 0.0f);
\end{lstlisting}

\gsub{Çapraz Çarpım (Cross Product)}

İki vektörün çapraz çarpımı, \textbf{her ikisine de dik} yeni bir vektördür.
Yüzey normallerini ve kamera baz eksenlerini bulmakta vazgeçilmezdir.
\[
\mathbf{a}\times\mathbf{b}=(a_yb_z-a_zb_y,\; a_zb_x-a_xb_z,\; a_xb_y-a_yb_x)
\]

\begin{lstlisting}
glm::vec3 n = glm::cross(a, b);   // vector perpendicular to both a and b

// Surface normal of a triangle (from two edges):
glm::vec3 edge1  = p1 - p0;
glm::vec3 edge2  = p2 - p0;
glm::vec3 normal = glm::normalize(glm::cross(edge1, edge2));
\end{lstlisting}

\bilgi{Çapraz çarpımın yönü \emph{sağ el kuralına} uyar ve sıra önemlidir:
$\mathbf{a}\times\mathbf{b}=-(\mathbf{b}\times\mathbf{a})$. Yüzey normalinin
"dışarı" mı "içeri" mi baktığı, köşe sıralamana (winding order) bağlıdır ---
bu, arka yüz eleme (back-face culling) için kritiktir.}

% =====================================================================
\gsec{Matrisler ve Dönüşümler}
% =====================================================================

Matris, bir sayı ızgarasıdır ve grafikte \textbf{dönüşümleri} temsil eder. Bir
noktayı bir matrisle çarpmak onu öteler, döndürür veya ölçekler. 3B grafikte
$4\times4$ matrisler kullanırız (nedenini homojen koordinatlarda göreceğiz).

\gsub{Matris Çarpımı ve Birim Matris}

Bir dönüşümü bir vektöre uygulamak, matris$\times$vektör çarpımıdır. Birden fazla
dönüşümü birleştirmek ise matris$\times$matris çarpımıdır. \textbf{Birim matris}
$\mathbf{I}$ hiçbir şeyi değiştirmez (1 ile çarpmak gibi):
\[
\mathbf{I}=\begin{pmatrix}1&0&0&0\\0&1&0&0\\0&0&1&0\\0&0&0&1\end{pmatrix}
\]

\ornek[Sayısal örnek: $2\times2$ matris $\times$ vektör]{Bir noktayı 90° döndüren
matris ile $(1,0)$ noktasını çarpalım:
\[
\begin{pmatrix}\cos90^\circ&-\sin90^\circ\\\sin90^\circ&\cos90^\circ\end{pmatrix}
\begin{pmatrix}1\\0\end{pmatrix}=
\begin{pmatrix}0&-1\\1&0\end{pmatrix}\begin{pmatrix}1\\0\end{pmatrix}=
\begin{pmatrix}0\\1\end{pmatrix}
\]
$(1,0)$ noktası $(0,1)$'e döndü --- yani saat yönünün tersine 90°. İşte tüm
dönüşümler bu satır$\times$sütun çarpımından ibarettir.}

\begin{lstlisting}
#include <glm/gtc/matrix_transform.hpp>

glm::mat4 identity(1.0f);  // constructing with 1.0f => IDENTITY (very common gotcha!)
\end{lstlisting}

\uyari{\code{glm::mat4 m;} demek matrisi \emph{sıfırlamaz} ve birim yapmaz;
ilklendirilmemiş kalabilir. Her zaman \code{glm::mat4 m(1.0f)} yazarak açıkça
birim matris kur. Sıfır matrisle çarpılan her nokta orijine çöker ve ekranda
hiçbir şey görünmez.}

\uyari{Matris çarpımı \textbf{değişmeli değildir}: $\mathbf{A}\mathbf{B}\neq
\mathbf{B}\mathbf{A}$. "Önce döndür sonra öte" ile "önce öte sonra döndür"
tamamen farklı sonuç verir. Sıra her şeydir.}

\gsub{Öteleme (Translation)}

Bir noktayı $(t_x, t_y, t_z)$ kadar kaydırır. Homojen $4\times4$ formu:
\[
\mathbf{T}=\begin{pmatrix}1&0&0&t_x\\0&1&0&t_y\\0&0&1&t_z\\0&0&0&1\end{pmatrix}
\]

\begin{lstlisting}
glm::mat4 translation =
    glm::translate(glm::mat4(1.0f), glm::vec3(2.0f, 1.0f, 0.0f));
// Moves the object +2 on X and +1 on Y.
\end{lstlisting}

\gsub{Ölçekleme (Scale)}

Her ekseni bağımsız büyütür/küçültür:
\[
\mathbf{S}=\begin{pmatrix}s_x&0&0&0\\0&s_y&0&0\\0&0&s_z&0\\0&0&0&1\end{pmatrix}
\]

\begin{lstlisting}
glm::mat4 scaling = glm::scale(glm::mat4(1.0f), glm::vec3(2.0f)); // 2x on every axis
\end{lstlisting}

\gsub{Döndürme (Rotation)}

Bir eksen etrafında $\theta$ açısıyla döndürür. Örneğin Z ekseni etrafında:
\[
\mathbf{R}_z(\theta)=\begin{pmatrix}\cos\theta&-\sin\theta&0&0\\\sin\theta&\cos\theta&0&0\\0&0&1&0\\0&0&0&1\end{pmatrix}
\]

\begin{lstlisting}
// glm::radians converts degrees -> radians (GLM expects radians!)
glm::mat4 rotation = glm::rotate(glm::mat4(1.0f),
                                 glm::radians(45.0f),          // angle
                                 glm::vec3(0.0f, 0.0f, 1.0f)); // around Z axis
\end{lstlisting}

\uyari{GLM'nin trigonometrik ve döndürme fonksiyonları \textbf{radyan} bekler,
derece değil. \code{glm::rotate(m, 90.0f, axis)} yazarsan 90 \emph{radyan}
dönersin (yani ~5157 derece). Daima \code{glm::radians(90.0f)} kullan.}

\gsub{Dönüşümleri Birleştirme}

Birden fazla dönüşümü tek matriste toplayabiliriz. Kritik nokta: GLM'de (ve
OpenGL'de) çarpım \textbf{sağdan sola} uygulanır. Yani kodda önce yazdığın matris,
noktaya en \emph{son} uygulanır.

\begin{lstlisting}
// Goal: first scale the object, then rotate it, and finally translate it.
// Application order to a vertex v: T * R * S * v  (S applies to v first)
glm::mat4 model(1.0f);
model = glm::translate(model, position);              // applied 3rd
model = glm::rotate(model, glm::radians(angle), axis); // applied 2nd
model = glm::scale(model, scale);                     // applied 1st (closest to v)
// Result: model = T * R * S
\end{lstlisting}

\notk[Sağdan sola sezgisi]{\code{model = T; model *= R; model *= S;} zinciri
$\mathbf{T}\cdot\mathbf{R}\cdot\mathbf{S}$ üretir. Bir köşe $\mathbf{v}$ bununla
çarpılınca önce $\mathbf{S}$ (ölçek), sonra $\mathbf{R}$ (döndür), en son
$\mathbf{T}$ (öte) etkir. Doğru sıra: önce yerel şekli düzenle (ölçek/döndür),
en son dünyaya yerleştir (öte). Tersini yaparsan nesne orijin etrafında yörüngeye
girer.}

% =====================================================================
\gsec{Homojen Koordinatlar}
% =====================================================================

Neden $3\times3$ değil de $4\times4$ matris? Çünkü \textbf{öteleme} bir matris
çarpımıyla ifade edilemez --- düz matris çarpımı orijini sabit tutar. Çözüm:
noktalara dördüncü bir bileşen $w$ ekleyip $(x, y, z, w)$ yapmak. Buna
\emph{homojen koordinatlar} denir.

\begin{center}
\begin{tabularx}{\textwidth}{@{}>{\bfseries\color{anaMavi}}l >{\raggedright\arraybackslash}X@{}}
\toprule
\rowcolor{acikMavi}\color{anaMavi}\textbf{$w$ değeri} & \color{anaMavi}\textbf{Anlamı}\\
\midrule
$w = 1$ & Bir \textbf{nokta} (konum). Öteleme onu etkiler.\\
\rowcolor{acikGri} $w = 0$ & Bir \textbf{yön} (vektör). Öteleme onu etkilemez, sadece döndürme/ölçek.\\
\bottomrule
\end{tabularx}
\end{center}

\begin{lstlisting}
glm::vec4 point     = glm::vec4(pos, 1.0f);  // w=1: a position, translation affects it
glm::vec4 direction = glm::vec4(dir, 0.0f);  // w=0: a direction, translation ignored

glm::mat4 model = /* ... */;
glm::vec3 worldPos = glm::vec3(model * point);      // translated + rotated + scaled
glm::vec3 worldDir = glm::vec3(model * direction);  // only rotated + scaled
\end{lstlisting}

\bilgi{Bu ayrım çok işe yarar: bir konum noktasını dünyaya taşırken (w=1) öteleme
uygulanmalı, ama bir yön vektörünü (ör. ışık yönü, yüzey normali, w=0) taşırken
öteleme \emph{uygulanmamalı} --- yönün başladığı yer önemsizdir.}

\gsub{Perspektif Bölme}

Perspektif projeksiyondan sonra $w$ artık 1 olmayabilir. GPU, clip uzayından
NDC'ye geçerken her bileşeni $w$'ye böler:
\[
(x, y, z, w) \;\longrightarrow\; \left(\tfrac{x}{w}, \tfrac{y}{w}, \tfrac{z}{w}\right)
\]
İşte "uzaktakiler küçük görünür" etkisi buradan gelir: uzak nesnelerde $w$
büyüktür, bölünce koordinatlar küçülür. Bu bölmeyi GPU otomatik yapar; biz
sadece doğru projeksiyon matrisini vermekle yükümlüyüz.

% =====================================================================
\gsec{View ve Projection: Kamera Matematiği}
% =====================================================================

\gsub{View Matrisi (LookAt)}

Aslında "kamerayı hareket ettirmek" diye bir şey yoktur; bunun yerine \emph{tüm
dünyayı} kameranın tersine hareket ettiririz. View matrisi bunu yapar. GLM'nin
\code{lookAt} fonksiyonu üç şeyle bir kamera kurar:

\begin{center}
\begin{tabularx}{\textwidth}{@{}>{\bfseries\color{anaMavi}}l >{\raggedright\arraybackslash}X@{}}
\toprule
\rowcolor{acikMavi}\color{anaMavi}\textbf{Parametre} & \color{anaMavi}\textbf{Anlamı}\\
\midrule
eye (position) & Kameranın dünyadaki konumu.\\
\rowcolor{acikGri} center (target) & Kameranın baktığı nokta.\\
up & "Yukarı" yönü (genelde $(0,1,0)$).\\
\bottomrule
\end{tabularx}
\end{center}

\begin{lstlisting}
glm::vec3 cameraPos    = glm::vec3(0.0f, 0.0f,  3.0f);
glm::vec3 cameraTarget = glm::vec3(0.0f, 0.0f,  0.0f);
glm::vec3 worldUp      = glm::vec3(0.0f, 1.0f,  0.0f);

glm::mat4 view = glm::lookAt(cameraPos, cameraTarget, worldUp);
\end{lstlisting}

\notk[Perde arkası]{\code{lookAt} içeride kamera için üç dik birim eksen kurar:
\emph{forward} (target$-$eye), \emph{right} (up $\times$ forward) ve \emph{true up}
(forward $\times$ right). Bu çapraz çarpımlar tam da 6. bölümde öğrendiğimiz işlemdir.
Sonra bu eksenlerle dünyayı kamera uzayına döndürüp öteler.}

\gsub{Projection Matrisi}

Projeksiyon matrisi 3B sahneyi 2B ekrana nasıl "yansıtacağımızı" belirler ve
görünür hacmi (frustum) tanımlar. İki tür vardır:

\begin{center}
\begin{tabularx}{\textwidth}{@{}>{\bfseries\color{anaMavi}}l >{\raggedright\arraybackslash}X@{}}
\toprule
\rowcolor{acikMavi}\color{anaMavi}\textbf{Perspektif} & \color{anaMavi}\textbf{Ortografik}\\
\midrule
Uzaktakiler küçülür (gerçekçi, 3B oyunlar). & Boyut mesafeden bağımsız (CAD, 2B, izometrik).\\
\rowcolor{acikGri} Görünür hacim bir \emph{piramit kesiği} (frustum). & Görünür hacim bir \emph{dikdörtgen prizma}.\\
\bottomrule
\end{tabularx}
\end{center}

\begin{lstlisting}
// Perspective: FOV, aspect ratio, near and far planes
float aspect = (float)width / (float)height;
glm::mat4 projection = glm::perspective(glm::radians(45.0f), // vertical field of view
                                        aspect,              // width / height
                                        0.1f,                // near (must be > 0!)
                                        100.0f);             // far

// Orthographic: left, right, bottom, top, near, far
glm::mat4 ortho = glm::ortho(0.0f, 800.0f, 0.0f, 600.0f, -1.0f, 1.0f);
\end{lstlisting}

\uyari{Perspektifte \code{near} düzlemi kesinlikle \textbf{0'dan büyük} olmalıdır
(ör. 0.1). 0 verirsen derinlik tamponu (z-buffer) hassasiyeti bozulur ve
"z-fighting" (titreşen yüzeyler) yaşarsın. Ayrıca near/far aralığını gereksiz
geniş tutma; hassasiyet near'a yakın toplanır.}

\gsub{Hepsini Birleştirmek}

Artık MVP zincirini tamamlayabiliriz. Vertex shader'da her köşeyi şöyle
dönüştüreceğiz:

\begin{lstlisting}
// On the C++ side we build three matrices and send them as uniforms:
glm::mat4 model      = /* object's placement in the world */;
glm::mat4 view       = glm::lookAt(cameraPos, cameraTarget, worldUp);
glm::mat4 projection = glm::perspective(glm::radians(45.0f), aspect, 0.1f, 100.0f);

// Inside the GLSL vertex shader:
//   gl_Position = projection * view * model * vec4(aPos, 1.0);
\end{lstlisting}

\ipucu{Kısım 2 bitti. Artık bir noktanın modelden ekrana yolculuğunun her
adımının matematiğini biliyorsun. Bundan sonrası bu matematiği çalışan bir
pencerede canlandırmak: GLFW ile pencere açıp OpenGL ile çizmek.}

% #####################################################################
\gpart{3}{GLFW: Pencere, Bağlam ve Girdi}
% #####################################################################

GLFW (Graphics Library Framework), çok platformlu, hafif bir C kütüphanesidir.
Tek işi vardır ama onu çok iyi yapar: bir OpenGL \textbf{bağlamı (context)} olan
pencere açmak, olayları (klavye, fare, boyutlandırma) yönetmek ve tamponları
değiştirmek. GLFW hiçbir şey \emph{çizmez} --- çizim OpenGL'in işidir. GLFW ise
üzerine çizeceğin tuvali ve girdi sistemini sağlar.

% =====================================================================
\gsec{Kurulum ve Derleme}
% =====================================================================

Bir OpenGL projesi tipik olarak üç parçadan oluşur:

\begin{center}
\begin{tabularx}{\textwidth}{@{}>{\bfseries\color{anaMavi}}l >{\raggedright\arraybackslash}X@{}}
\toprule
\rowcolor{acikMavi}\color{anaMavi}\textbf{Bileşen} & \color{anaMavi}\textbf{Görevi}\\
\midrule
GLFW & Pencere + bağlam + girdi.\\
\rowcolor{acikGri} GLAD (veya GLEW) & OpenGL fonksiyon işaretçilerini çalışma anında yükler.\\
GLM & Vektör/matris matematiği (başlık-only).\\
\bottomrule
\end{tabularx}
\end{center}

\bilgi{\textbf{GLAD neden gerekli?} OpenGL'in modern fonksiyonları (ör.
\code{glCreateShader}) sürücüde bulunur ve adresleri çalışma anında öğrenilmelidir.
GLAD bu adresleri senin için yükler. GLAD kodunu \url{https://glad.dav1d.de}
sitesinden üretirsin: API=OpenGL, Version=3.3, Profile=Core seçip indir.}

\gsub{Linux'ta Paketlerin Kurulumu}

\begin{lstlisting}[style=shstyle]
# Debian / Ubuntu
$ sudo apt install libglfw3-dev libglm-dev cmake g++
# Download GLAD from glad.dav1d.de and add it (glad.c + include/) to the project

# Fedora
$ sudo dnf install glfw-devel glm-devel cmake gcc-c++

# Arch
$ sudo pacman -S glfw glm cmake
\end{lstlisting}

\gsub{CMake ile Proje Yapısı}

Önerilen dizin düzeni:

\begin{lstlisting}[style=shstyle]
project/
|-- CMakeLists.txt
|-- src/
|   \-- main.cpp
|-- glad/
|   |-- src/glad.c
|   |-- include/glad/glad.h
|   \-- include/KHR/khrplatform.h
\-- shaders/
    |-- basic.vert
    \-- basic.frag
\end{lstlisting}

\begin{lstlisting}[style=shstyle]
# CMakeLists.txt
cmake_minimum_required(VERSION 3.16)
project(GlfwGuide LANGUAGES C CXX)

set(CMAKE_CXX_STANDARD 17)
set(CMAKE_CXX_STANDARD_REQUIRED ON)

find_package(glfw3 REQUIRED)
find_package(OpenGL REQUIRED)

add_executable(app
    src/main.cpp
    glad/src/glad.c)          # include GLAD in the build

target_include_directories(app PRIVATE glad/include)
target_link_libraries(app PRIVATE glfw OpenGL::GL ${CMAKE_DL_LIBS})
# ${CMAKE_DL_LIBS}: GLAD may need libdl for dlopen
\end{lstlisting}

\begin{lstlisting}[style=shstyle]
# Build
$ cmake -S . -B build
$ cmake --build build
$ ./build/app
\end{lstlisting}

\notk[Elle derleme]{CMake istemiyorsan tek satırla da derleyebilirsin:
\code{g++ -std=c++17 src/main.cpp glad/src/glad.c -Iglad/include -lglfw -lGL -ldl -o app}}

% =====================================================================
\gsec{İlk Pencere}
% =====================================================================

Şimdi en küçük tam GLFW programını yazalım. Bu program bir pencere açar, koyu
turkuaz bir arka planla temizler ve kullanıcı kapatana kadar açık tutar. Her
satırı yorumluyoruz.

\begin{lstlisting}
#include <glad/glad.h>   // GLAD must be included BEFORE GLFW!
#include <GLFW/glfw3.h>
#include <iostream>

// GLFW reports internal errors through this callback
void errorCallback(int code, const char* description) {
    std::cerr << "GLFW error " << code << ": " << description << '\n';
}

int main() {
    // 1) Register the error callback (before glfwInit so init errors are caught)
    glfwSetErrorCallback(errorCallback);

    // 2) Initialize GLFW
    if (!glfwInit()) {
        std::cerr << "Failed to initialize GLFW\n";
        return -1;
    }

    // 3) Request OpenGL 3.3 Core profile
    glfwWindowHint(GLFW_CONTEXT_VERSION_MAJOR, 3);
    glfwWindowHint(GLFW_CONTEXT_VERSION_MINOR, 3);
    glfwWindowHint(GLFW_OPENGL_PROFILE, GLFW_OPENGL_CORE_PROFILE);
#ifdef __APPLE__
    glfwWindowHint(GLFW_OPENGL_FORWARD_COMPAT, GL_TRUE); // required on macOS
#endif

    // 4) Create the window (width, height, title, monitor, share)
    GLFWwindow* window = glfwCreateWindow(800, 600, "First Window", nullptr, nullptr);
    if (!window) {
        std::cerr << "Failed to create window\n";
        glfwTerminate();
        return -1;
    }

    // 5) Make this window's OpenGL context current on the calling thread
    glfwMakeContextCurrent(window);

    // 6) Load OpenGL function pointers via GLAD (AFTER the context is current!)
    if (!gladLoadGLLoader((GLADloadproc)glfwGetProcAddress)) {
        std::cerr << "Failed to initialize GLAD\n";
        return -1;
    }

    // 7) Set the region OpenGL renders into
    glViewport(0, 0, 800, 600);

    // 8) MAIN LOOP: run until the window is asked to close
    while (!glfwWindowShouldClose(window)) {
        // Clear the background
        glClearColor(0.2f, 0.3f, 0.3f, 1.0f); // dark teal
        glClear(GL_COLOR_BUFFER_BIT);

        // Swap the front/back buffers (present the frame)
        glfwSwapBuffers(window);
        // Process pending events (keyboard, mouse, close)
        glfwPollEvents();
    }

    // 9) Release resources
    glfwDestroyWindow(window);
    glfwTerminate();
    return 0;
}
\end{lstlisting}

\uyari{\code{glad.h} her zaman \code{glfw3.h}'den \textbf{önce} dahil edilmelidir,
çünkü GLAD OpenGL başlıklarını kendi tanımlar ve GLFW bunu algılar. Sıra tersse
"gl.h already included" hataları alırsın. Ayrıca GLAD'i \emph{ancak}
\code{glfwMakeContextCurrent}'tan sonra yükle --- fonksiyonlar bir bağlam
olmadan çözülemez.}

\gsub{Ana Döngünün Anatomisi}

Her gerçek zamanlı uygulamanın kalbi bu döngüdür. Her tur bir \emph{kare}dir:

\begin{center}
\begin{tabularx}{\textwidth}{@{}>{\bfseries\color{anaMavi}}l >{\raggedright\arraybackslash}X@{}}
\toprule
\rowcolor{acikMavi}\color{anaMavi}\textbf{Adım} & \color{anaMavi}\textbf{Çağrı}\\
\midrule
Kapatma kontrolü & \code{glfwWindowShouldClose(window)}\\
\rowcolor{acikGri} Girdi işle & \code{processInput()} (bizim fonksiyon)\\
Ekranı temizle & \code{glClear(...)}\\
\rowcolor{acikGri} Çizim & \code{glDrawArrays / glDrawElements}\\
Tamponları değiştir & \code{glfwSwapBuffers(window)}\\
\rowcolor{acikGri} Olayları işle & \code{glfwPollEvents()}\\
\bottomrule
\end{tabularx}
\end{center}

% =====================================================================
\gsec{Callback'ler ve Olay İşleme}
% =====================================================================

GLFW olayları (pencere boyutu değişti, tuşa basıldı, fare oynadı) \textbf{callback}
fonksiyonlarıyla bildirir. Sen bir fonksiyon yazar, onu GLFW'ye kaydedersin; GLFW
uygun olay olduğunda onu çağırır.

\gsub{Pencere Boyutlandırma}

Kullanıcı pencereyi yeniden boyutlandırınca viewport'u güncellemezsen çizimin
gerilir. Bunu bir framebuffer-size callback ile çözeriz:

\begin{lstlisting}
// GLFW calls this whenever the framebuffer (drawable area) is resized
void framebufferSizeCallback(GLFWwindow* window, int width, int height) {
    glViewport(0, 0, width, height); // update OpenGL's drawing region
}

int main() {
    // ... window setup ...
    glfwSetFramebufferSizeCallback(window, framebufferSizeCallback);
    // ...
}
\end{lstlisting}

\notk[Neden framebuffer boyutu?]{Yüksek-DPI (Retina) ekranlarda pencere boyutu
(ekran birimi) ile framebuffer boyutu (gerçek piksel) farklı olabilir.
\code{glViewport} gerçek pikselleri ister, o yüzden \emph{framebuffer} boyutu
callback'ini kullanırız, pencere boyutunu değil.}

\gsub{Klavye Girdisi: İki Yöntem}

Klavyeyi iki şekilde okuyabilirsin. Sürekli hareket için \emph{polling}, tek
seferlik olaylar için \emph{callback} uygundur.

\begin{lstlisting}
// METHOD 1: Polling — query state every frame (ideal for continuous movement)
void processInput(GLFWwindow* window) {
    if (glfwGetKey(window, GLFW_KEY_ESCAPE) == GLFW_PRESS)
        glfwSetWindowShouldClose(window, true); // quit on ESC

    if (glfwGetKey(window, GLFW_KEY_W) == GLFW_PRESS)
        std::cout << "Moving forward...\n";
}

// METHOD 2: Callback — event based (ideal for single presses, menus, shortcuts)
void keyCallback(GLFWwindow* win, int key, int scancode, int action, int mods) {
    if (key == GLFW_KEY_SPACE && action == GLFW_PRESS)
        std::cout << "Fire!\n";                  // once, at the moment of press
    if (key == GLFW_KEY_F && action == GLFW_PRESS && (mods & GLFW_MOD_CONTROL))
        std::cout << "Ctrl+F pressed\n";         // modifier keys
}

int main() {
    // ...
    glfwSetKeyCallback(window, keyCallback);
    while (!glfwWindowShouldClose(window)) {
        processInput(window);  // polling happens here
        // ... draw ...
    }
}
\end{lstlisting}

\begin{center}
\begin{tabularx}{\textwidth}{@{}>{\bfseries\color{anaMavi}}l >{\raggedright\arraybackslash}X@{}}
\toprule
\rowcolor{acikMavi}\color{anaMavi}\textbf{action değeri} & \color{anaMavi}\textbf{Anlamı}\\
\midrule
GLFW\_PRESS & Tuşa yeni basıldı.\\
\rowcolor{acikGri} GLFW\_RELEASE & Tuş bırakıldı.\\
GLFW\_REPEAT & Tuş basılı tutuluyor (tekrar).\\
\bottomrule
\end{tabularx}
\end{center}

\ipucu{Karakter/oyuncu hareketi için \textbf{polling} kullan (\code{glfwGetKey}):
her karede "W basılı mı?" diye sor, basılıysa hareket ettir. Menü seçimi, zıplama,
ateş gibi \emph{tek seferlik} eylemler için \textbf{callback} kullan --- böylece
tuş basılı tutulunca eylem tekrarlanmaz.}

\gsub{Fare Girdisi}

Fare konumu, tuşları ve tekerleği için ayrı callback'ler vardır. FPS kamera için
fare hareketini kullanacağız (Kısım 4).

\begin{lstlisting}
// Mouse movement (cursor position, in screen pixels)
void cursorPosCallback(GLFWwindow* window, double xpos, double ypos) {
    std::cout << "Cursor: " << xpos << ", " << ypos << "\n";
}

// Scroll wheel (horizontal and vertical offsets)
void scrollCallback(GLFWwindow* window, double xoffset, double yoffset) {
    std::cout << "Scroll: " << yoffset << "\n"; // ideal for zoom
}

// Mouse buttons
void mouseButtonCallback(GLFWwindow* win, int button, int action, int mods) {
    if (button == GLFW_MOUSE_BUTTON_LEFT && action == GLFW_PRESS)
        std::cout << "Left click\n";
}

int main() {
    // ...
    glfwSetCursorPosCallback(window, cursorPosCallback);
    glfwSetScrollCallback(window, scrollCallback);
    glfwSetMouseButtonCallback(window, mouseButtonCallback);

    // For an FPS camera: hide the cursor and lock it to the window
    glfwSetInputMode(window, GLFW_CURSOR, GLFW_CURSOR_DISABLED);
}
\end{lstlisting}

% =====================================================================
\gsec{Zaman, Delta Time ve FPS}
% =====================================================================

Farklı bilgisayarlar farklı hızlarda kare üretir. Hareketi kare sayısına
bağlarsan, hızlı makinede nesne uçar, yavaşta sürünür. Çözüm: her hareketi
\textbf{delta time} (iki kare arası geçen süre) ile çarpmak. Böylece hız
donanımdan bağımsız olur.

\begin{lstlisting}
float deltaTime = 0.0f;  // duration of the last frame (seconds)
float lastFrame = 0.0f;  // timestamp of the previous frame

int main() {
    // ...
    while (!glfwWindowShouldClose(window)) {
        float currentFrame = (float)glfwGetTime(); // seconds since program start
        deltaTime = currentFrame - lastFrame;
        lastFrame = currentFrame;

        // Scale movement by delta time (units: distance per second)
        float velocity = 2.5f * deltaTime;
        cameraPos += velocity * direction;   // now independent of FPS!

        // ... draw ...
        glfwSwapBuffers(window);
        glfwPollEvents();
    }
}
\end{lstlisting}

\gsub{Basit FPS Sayacı}

\begin{lstlisting}
double previousSeconds = glfwGetTime();
int    frameCount = 0;

// Inside the main loop:
double currentSeconds = glfwGetTime();
frameCount++;
if (currentSeconds - previousSeconds >= 1.0) {  // one second elapsed
    std::cout << "FPS: " << frameCount << '\n';
    frameCount = 0;
    previousSeconds = currentSeconds;
}
\end{lstlisting}

\ipucu{\code{glfwSwapInterval(1)} ile VSync açarak FPS'i ekran yenileme hızına
sabitleyebilir, GPU'yu boşuna yormazsın. Performans ölçerken \code{glfwSwapInterval(0)}
ile kapat ki gerçek kare süresini görebilesin.}

% #####################################################################
\gpart{4}{OpenGL ile Çizim}
% #####################################################################

Artık pencere ve girdi hazır. Bu kısımda gerçek çizime geçiyoruz: köşe verisini
GPU'ya yükleyip, shader'larla işleyip, ekrana üçgen, kare, dokulu yüzey ve nihayet
aydınlatılmış 3B küp çizeceğiz. Kısım 2'nin matematiği burada canlanacak.

% =====================================================================
\gsec{Shader'lar ve GLSL Dili Derinlemesine}
% =====================================================================

Shader, GPU üzerinde çalışan küçük bir programdır. Modern OpenGL'de shader olmadan
tek bir piksel bile çizemezsin --- bu yüzden shader'ları gerçekten anlamak
zorunludur. Shader'lar \textbf{GLSL} (OpenGL Shading Language) ile yazılır; C'ye
çok benzeyen ama grafik için özelleştirilmiş bir dildir.

\gsub{Shader Türleri ve Boru Hattındaki Yeri}

\begin{center}
\begin{tabularx}{\textwidth}{@{}>{\bfseries\color{anaMavi}}l c >{\raggedright\arraybackslash}X@{}}
\toprule
\rowcolor{acikMavi}\color{anaMavi}\textbf{Shader} & \color{anaMavi}\textbf{Zorunlu?} & \color{anaMavi}\textbf{Ne yapar?}\\
\midrule
Vertex & Evet & Her \emph{köşe} için bir kez çalışır; konumu clip uzayına koyar, öznitelikleri sonraki aşamaya geçirir.\\
\rowcolor{acikGri} Fragment & Evet & Her \emph{fragment} (aday piksel) için bir kez çalışır; nihai rengi üretir.\\
Geometry & Hayır & Primitive'leri işler; yeni geometri üretebilir (ör. tek noktadan quad).\\
\rowcolor{acikGri} Tessellation & Hayır & Yüzeyleri dinamik olarak alt bölümlere ayırır (arazi, LOD).\\
Compute & Hayır & Grafikten bağımsız genel amaçlı GPU hesabı (parçacık, fizik).\\
\bottomrule
\end{tabularx}
\end{center}

\bilgi{Bu rehberde yalnızca zorunlu ikiliyi --- \textbf{vertex} ve \textbf{fragment}
--- kullanıyoruz. İkisi bir \emph{shader programı} altında derlenip bağlanır.
Vertex'in \code{out} değişkenleri, fragment'in \code{in} değişkenlerine
\emph{isim üzerinden} bağlanır ve arada rasterizasyon interpolasyon yapar.}

\gsub{Bir Shader'ın İskeleti}

\begin{lstlisting}[style=glslstyle]
#version 330 core   // GLSL version: OpenGL 3.3 Core -> GLSL 330

// (input/output/uniform declarations go here)

void main() {
    // shader body — runs once per vertex/fragment
}
\end{lstlisting}

\notk[Sürüm eşleşmesi]{\code{\#version 330 core}, OpenGL 3.3 Core bağlamıyla
eşleşir. GLSL sürüm numarası OpenGL'inkinden farklıdır: GL 3.3 $\to$ GLSL 330,
GL 4.6 $\to$ GLSL 460. \code{\#version} satırı dosyanın \textbf{ilk} satırı
olmalıdır --- önünde yorum bile olamaz.}

\gsub{Veri Tipleri}

\begin{center}
\begin{tabularx}{\textwidth}{@{}>{\ttfamily\bfseries\color{vurguMavi}}l >{\raggedright\arraybackslash}X@{}}
\toprule
\rowcolor{acikMavi}\color{anaMavi}\textbf{Tip} & \color{anaMavi}\textbf{Açıklama}\\
\midrule
float, int, bool & Skaler tipler.\\
\rowcolor{acikGri} vec2, vec3, vec4 & 2/3/4 bileşenli float vektör (konum, renk, UV).\\
ivec3, bvec3 & int / bool vektörleri.\\
\rowcolor{acikGri} mat2, mat3, mat4 & Kare matrisler (dönüşümler).\\
sampler2D & Bir dokuya erişim tutamacı.\\
\bottomrule
\end{tabularx}
\end{center}

\begin{lstlisting}[style=glslstyle]
vec3 color    = vec3(1.0, 0.5, 0.2);       // orange
vec4 position = vec4(color, 1.0);          // vec3 + w -> vec4 (constructor stitching)
vec3 uniform3 = vec3(0.5);                 // (0.5, 0.5, 0.5) — all the same
mat4 identity = mat4(1.0);                 // identity matrix
\end{lstlisting}

\gsub{Swizzling: Vektör Bileşenlerine Erişim}

GLSL'in en sevilen özelliği \textbf{swizzling}'dir: vektör bileşenlerini harflerle
seçip yeniden düzenleyebilirsin. \code{x,y,z,w} (konum), \code{r,g,b,a} (renk) ve
\code{s,t,p,q} (doku) aynı bileşenlere farklı isimlerdir.

\begin{lstlisting}[style=glslstyle]
vec4 v = vec4(1.0, 2.0, 3.0, 4.0);

float x   = v.x;         // 1.0
vec2  xy  = v.xy;        // (1.0, 2.0)
vec3  bgr = v.bgr;       // (3.0, 2.0, 1.0) — reverse color channels
vec3  all = v.xxx;       // (1.0, 1.0, 1.0) — repetition is allowed
v.xy = v.yx;             // swap components
\end{lstlisting}

\ipucu{Swizzling sadece okuma değil, atama için de çalışır. \code{fragColor.rgb *=
0.5;} tüm renk kanallarını yarıya indirirken alfayı korur.}

\gsub{Giriş, Çıkış ve Nitelikler (in / out / uniform)}

\begin{center}
\begin{tabularx}{\textwidth}{@{}>{\ttfamily\bfseries\color{turuncu}}l >{\raggedright\arraybackslash}X@{}}
\toprule
\rowcolor{acikMavi}\color{anaMavi}\textbf{Nitelik} & \color{anaMavi}\textbf{Anlamı}\\
\midrule
in & Bu shader'a \emph{gelen} veri. Vertex'te: VBO'dan öznitelik. Fragment'te: vertex'ten interpolasyonla gelen değer.\\
\rowcolor{acikGri} out & Bu shader'dan \emph{çıkan} veri. Vertex'te: sonraki aşamaya. Fragment'te: framebuffer'a (renk).\\
uniform & CPU'dan (C++) gelen, tüm çizim boyunca \emph{sabit} global değer.\\
\rowcolor{acikGri} layout(location=N) & Bir \code{in} özniteliğine sabit numara atar; \code{glVertexAttribPointer}'daki index ile eşleşir.\\
\bottomrule
\end{tabularx}
\end{center}

\begin{lstlisting}[style=glslstyle]
// VERTEX SHADER — example showing all three qualifiers
#version 330 core
layout (location = 0) in vec3 aPos;     // position attribute from the VBO
layout (location = 1) in vec3 aColor;   // color attribute from the VBO

out vec3 vColor;        // passed to the fragment shader (gets interpolated)
uniform mat4 mvp;       // constant transform matrix from the CPU

void main() {
    gl_Position = mvp * vec4(aPos, 1.0);   // required output: clip-space position
    vColor = aColor;                       // forward color to the next stage
}
\end{lstlisting}

\begin{lstlisting}[style=glslstyle]
// FRAGMENT SHADER — receives the vertex shader's out as in
#version 330 core
in  vec3 vColor;        // interpolated color from the vertex shader
out vec4 FragColor;     // final color written to the framebuffer
uniform float alpha;    // transparency from the CPU

void main() {
    FragColor = vec4(vColor, alpha);
}
\end{lstlisting}

\uyari{Vertex'teki \code{out} ile fragment'teki \code{in} \textbf{aynı ada ve tipe}
sahip olmalıdır --- eşleşme isim üzerindendir. Adı yanlış yazarsan fragment o
değeri hiç almaz ve genelde siyah ekranla karşılaşırsın.}

\gsub{Yerleşik Değişkenler ve Fonksiyonlar}

\begin{center}
\begin{tabularx}{\textwidth}{@{}>{\ttfamily\bfseries\color{turuncu}}l c >{\raggedright\arraybackslash}X@{}}
\toprule
\rowcolor{acikMavi}\color{anaMavi}\textbf{Değişken} & \color{anaMavi}\textbf{Shader} & \color{anaMavi}\textbf{Rolü}\\
\midrule
gl\_Position & Vertex & \textbf{Zorunlu çıktı}: köşenin clip uzayı konumu (vec4).\\
\rowcolor{acikGri} gl\_FragCoord & Fragment & Fragment'in ekran piksel koordinatı (salt-okunur).\\
gl\_VertexID & Vertex & O anki köşenin indeksi.\\
\bottomrule
\end{tabularx}
\end{center}

\begin{lstlisting}[style=glslstyle]
// A helper function + built-ins in a fragment shader
float luminance(vec3 c) {
    return dot(c, vec3(0.2126, 0.7152, 0.0722)); // grayscale weights
}

void main() {
    vec3 texColor = texture(image, uv).rgb;
    float gray    = luminance(texColor);
    vec3 result   = mix(texColor, vec3(gray), 0.5); // 50% desaturated
    FragColor = vec4(result, 1.0);
}
\end{lstlisting}

\bilgi{Fragment shader'da \code{discard;} anahtar kelimesi o fragment'i tamamen
atar (hiç yazılmaz) --- ağaç yaprağı gibi delikli dokularda saydam pikselleri
elemek için kullanılır. GLSL'de \code{if}/\code{for} vardır ama GPU'da dallanma
pahalı olabilir; mümkünse \code{mix}/\code{step} gibi dalsız işlevleri tercih et.}

\ipucu{Shader'ları anlamanın en hızlı yolu \emph{oynamaktır}: bir zaman uniform'u
alıp \code{sin} ile renkleri dalgalandır. \url{https://www.shadertoy.com} tam da
bu deneyler için bir oyun alanıdır.}

% =====================================================================
\gsec{VBO, VAO ve İlk Üçgen}
% =====================================================================

Modern OpenGL'de üçgen çizmek için üç kavramı bilmemiz gerekir:

\begin{center}
\begin{tabularx}{\textwidth}{@{}>{\bfseries\color{anaMavi}}l >{\raggedright\arraybackslash}X@{}}
\toprule
\rowcolor{acikMavi}\color{anaMavi}\textbf{Nesne} & \color{anaMavi}\textbf{Görevi}\\
\midrule
VBO (Vertex Buffer Object) & Köşe verisini (konum, renk...) GPU belleğine koyar.\\
\rowcolor{acikGri} VAO (Vertex Array Object) & Bu verinin \emph{nasıl yorumlanacağını} kaydeder (öznitelik düzeni).\\
Shader Program & Vertex + fragment shader'ın derlenmiş, bağlanmış hali.\\
\bottomrule
\end{tabularx}
\end{center}

\gsub{Shader'ları Yazmak}

\begin{lstlisting}[style=glslstyle]
// basic.vert — VERTEX SHADER
#version 330 core
layout (location = 0) in vec3 aPos;   // attribute 0: position

void main() {
    // Put the position straight into clip space (w=1). No matrix — vertices are in NDC.
    gl_Position = vec4(aPos.x, aPos.y, aPos.z, 1.0);
}
\end{lstlisting}

\begin{lstlisting}[style=glslstyle]
// basic.frag — FRAGMENT SHADER
#version 330 core
out vec4 FragColor;   // this fragment's final color

void main() {
    FragColor = vec4(1.0, 0.5, 0.2, 1.0);  // orange (RGBA)
}
\end{lstlisting}

\gsub{Shader Derleme ve Bağlama}

Shader kaynak kodu bir string'dir; onu GPU'da derleyip programa bağlamalıyız.
Hata kontrolü şart --- yoksa siyah ekranla baş başa kalırsın.

\begin{lstlisting}
unsigned int compileShader(unsigned int type, const char* source) {
    unsigned int shader = glCreateShader(type); // GL_VERTEX_SHADER / GL_FRAGMENT_SHADER
    glShaderSource(shader, 1, &source, nullptr);
    glCompileShader(shader);

    int success;
    glGetShaderiv(shader, GL_COMPILE_STATUS, &success);
    if (!success) {
        char log[512];
        glGetShaderInfoLog(shader, 512, nullptr, log);
        std::cerr << "Shader compile error:\n" << log << "\n";
    }
    return shader;
}

unsigned int createShaderProgram(const char* vertexSrc, const char* fragmentSrc) {
    unsigned int vs = compileShader(GL_VERTEX_SHADER,   vertexSrc);
    unsigned int fs = compileShader(GL_FRAGMENT_SHADER, fragmentSrc);

    unsigned int program = glCreateProgram();
    glAttachShader(program, vs);
    glAttachShader(program, fs);
    glLinkProgram(program);

    int success;
    glGetProgramiv(program, GL_LINK_STATUS, &success);
    if (!success) {
        char log[512];
        glGetProgramInfoLog(program, 512, nullptr, log);
        std::cerr << "Program link error:\n" << log << "\n";
    }
    // The compiled shaders are now embedded in the program; delete the singles
    glDeleteShader(vs);
    glDeleteShader(fs);
    return program;
}
\end{lstlisting}

\gsub{Köşe Verisi, VAO/VBO Kurulumu ve Çizim}

\begin{lstlisting}
// Three vertices in NDC: bottom-left, bottom-right, top-center
float vertices[] = {
    -0.5f, -0.5f, 0.0f,   // bottom left
     0.5f, -0.5f, 0.0f,   // bottom right
     0.0f,  0.5f, 0.0f    // top center
};

unsigned int VAO, VBO;
glGenVertexArrays(1, &VAO);
glGenBuffers(1, &VBO);

glBindVertexArray(VAO);                     // 1) bind the VAO (recording begins)

glBindBuffer(GL_ARRAY_BUFFER, VBO);         // 2) bind the VBO
glBufferData(GL_ARRAY_BUFFER, sizeof(vertices), vertices, GL_STATIC_DRAW); // 3) upload

// 4) Describe attribute 0: 3 floats, position; stride = 3 floats; offset = 0
glVertexAttribPointer(0, 3, GL_FLOAT, GL_FALSE, 3 * sizeof(float), (void*)0);
glEnableVertexAttribArray(0);               // 5) enable attribute 0

// (optional) unbind — the VAO already recorded the layout
glBindBuffer(GL_ARRAY_BUFFER, 0);
glBindVertexArray(0);
\end{lstlisting}

\begin{lstlisting}
// Inside the main loop — drawing:
glUseProgram(shaderProgram);   // use our shader
glBindVertexArray(VAO);        // restore the layout
glDrawArrays(GL_TRIANGLES, 0, 3); // start at 0, draw 3 vertices
\end{lstlisting}

\notk[glVertexAttribPointer'ın parametreleri]{Bu fonksiyon OpenGL'de en çok kafa
karıştırandır. \code{(index, size, type, normalized, stride, offset)}: \emph{index}
= shader'daki \code{location}; \emph{size} = bileşen sayısı (vec3 için 3);
\emph{stride} = bir köşenin toplam bayt boyu; \emph{offset} = bu özniteliğin köşe
içindeki başlangıç baytı.}

\uyari{Hiçbir şey görünmüyorsa kontrol listesi: (1) Shader'lar derlendi/bağlandı mı
(log'a bak)? (2) \code{glUseProgram} çağrıldı mı? (3) VAO bağlı mı? (4) Köşeler
NDC $[-1,1]$ içinde mi? (5) Üçgen sarım yönü yüzünden culling'e mi takıldı?
(6) \code{glClear} rengi ile çizim rengi aynı mı (görünmez üçgen)?}

\gsub{Tam, Kendi Kendine Yeten İlk Üçgen Programı}

Yukarıdaki parçaları tek, derlenip çalışan bir dosyada birleştirelim:

\begin{lstlisting}
#include <glad/glad.h>
#include <GLFW/glfw3.h>
#include <iostream>

const char* vertexSrc = R"(
#version 330 core
layout (location = 0) in vec3 aPos;
void main() { gl_Position = vec4(aPos, 1.0); }
)";

const char* fragmentSrc = R"(
#version 330 core
out vec4 FragColor;
void main() { FragColor = vec4(1.0, 0.5, 0.2, 1.0); }
)";

int main() {
    glfwInit();
    glfwWindowHint(GLFW_CONTEXT_VERSION_MAJOR, 3);
    glfwWindowHint(GLFW_CONTEXT_VERSION_MINOR, 3);
    glfwWindowHint(GLFW_OPENGL_PROFILE, GLFW_OPENGL_CORE_PROFILE);

    GLFWwindow* window = glfwCreateWindow(800, 600, "Triangle", nullptr, nullptr);
    glfwMakeContextCurrent(window);
    gladLoadGLLoader((GLADloadproc)glfwGetProcAddress);

    // Compile shaders (inline for brevity)
    unsigned int vs = glCreateShader(GL_VERTEX_SHADER);
    glShaderSource(vs, 1, &vertexSrc, nullptr);   glCompileShader(vs);
    unsigned int fs = glCreateShader(GL_FRAGMENT_SHADER);
    glShaderSource(fs, 1, &fragmentSrc, nullptr); glCompileShader(fs);
    unsigned int program = glCreateProgram();
    glAttachShader(program, vs); glAttachShader(program, fs);
    glLinkProgram(program);
    glDeleteShader(vs); glDeleteShader(fs);

    float vertices[] = {
        -0.5f, -0.5f, 0.0f,
         0.5f, -0.5f, 0.0f,
         0.0f,  0.5f, 0.0f
    };
    unsigned int VAO, VBO;
    glGenVertexArrays(1, &VAO);
    glGenBuffers(1, &VBO);
    glBindVertexArray(VAO);
    glBindBuffer(GL_ARRAY_BUFFER, VBO);
    glBufferData(GL_ARRAY_BUFFER, sizeof(vertices), vertices, GL_STATIC_DRAW);
    glVertexAttribPointer(0, 3, GL_FLOAT, GL_FALSE, 3 * sizeof(float), (void*)0);
    glEnableVertexAttribArray(0);

    while (!glfwWindowShouldClose(window)) {
        if (glfwGetKey(window, GLFW_KEY_ESCAPE) == GLFW_PRESS)
            glfwSetWindowShouldClose(window, true);

        glClearColor(0.2f, 0.3f, 0.3f, 1.0f);
        glClear(GL_COLOR_BUFFER_BIT);

        glUseProgram(program);
        glBindVertexArray(VAO);
        glDrawArrays(GL_TRIANGLES, 0, 3);

        glfwSwapBuffers(window);
        glfwPollEvents();
    }

    glDeleteVertexArrays(1, &VAO);
    glDeleteBuffers(1, &VBO);
    glDeleteProgram(program);
    glfwTerminate();
    return 0;
}
\end{lstlisting}

\ipucu{Bu tam programı derleyip çalıştırdığında koyu turkuaz zeminde turuncu bir
üçgen görürsün. OpenGL yolculuğunun "merhaba dünya"sıdır --- buradan sonrası
detay eklemekten ibaret.}

% =====================================================================
\gsec{Öznitelikler ve Renk İnterpolasyonu}
% =====================================================================

Şimdi her köşeye \emph{renk} de ekleyelim. Rasterizasyon, köşe renklerini üçgen
yüzeyi boyunca otomatik \textbf{interpolasyon} yapar --- ortadaki piksel üç rengin
karışımı olur.

\begin{lstlisting}
// Each vertex: 3 position + 3 color = 6 floats
float vertices[] = {
    // position          // color
    -0.5f, -0.5f, 0.0f,  1.0f, 0.0f, 0.0f,  // bottom left  — red
     0.5f, -0.5f, 0.0f,  0.0f, 1.0f, 0.0f,  // bottom right — green
     0.0f,  0.5f, 0.0f,  0.0f, 0.0f, 1.0f   // top          — blue
};

glBindVertexArray(VAO);
glBindBuffer(GL_ARRAY_BUFFER, VBO);
glBufferData(GL_ARRAY_BUFFER, sizeof(vertices), vertices, GL_STATIC_DRAW);

// Attribute 0: position (offset 0)
glVertexAttribPointer(0, 3, GL_FLOAT, GL_FALSE, 6 * sizeof(float), (void*)0);
glEnableVertexAttribArray(0);
// Attribute 1: color (offset 3 floats = 12 bytes)
glVertexAttribPointer(1, 3, GL_FLOAT, GL_FALSE, 6 * sizeof(float),
                      (void*)(3 * sizeof(float)));
glEnableVertexAttribArray(1);
\end{lstlisting}

\begin{lstlisting}[style=glslstyle]
// VERTEX SHADER — forward the color to the fragment shader
#version 330 core
layout (location = 0) in vec3 aPos;
layout (location = 1) in vec3 aColor;

out vec3 vColor;   // bridge to the fragment shader (gets interpolated)

void main() {
    gl_Position = vec4(aPos, 1.0);
    vColor = aColor;
}
\end{lstlisting}

\begin{lstlisting}[style=glslstyle]
// FRAGMENT SHADER — use the interpolated color
#version 330 core
in  vec3 vColor;    // comes from the vertex shader, SPREAD across the surface
out vec4 FragColor;

void main() {
    FragColor = vec4(vColor, 1.0);
}
\end{lstlisting}

\bilgi{Vertex shader'daki \code{out vec3 vColor}, fragment shader'daki \code{in vec3
vColor} ile \emph{isim üzerinden} eşleşir. GPU köşe değerlerini üçgen yüzeyinde
perspektif-doğru interpolasyonla yayar. Sonuç: köşeleri kırmızı/yeşil/mavi olan,
ortası yumuşak geçişli bir üçgen.}

% =====================================================================
\gsec{EBO ile Dikdörtgen ve Uniform'lar}
% =====================================================================

\gsub{EBO (Element Buffer Object)}

Bir dikdörtgen iki üçgenden oluşur (6 köşe), ama sadece 4 benzersiz köşe vardır.
Köşeleri tekrar etmemek için \textbf{EBO} kullanırız: köşeleri bir kez tanımlar,
indekslerle hangi köşenin hangi üçgende olduğunu belirtiriz.

\begin{lstlisting}
float vertices[] = {
     0.5f,  0.5f, 0.0f,  // 0: top right
     0.5f, -0.5f, 0.0f,  // 1: bottom right
    -0.5f, -0.5f, 0.0f,  // 2: bottom left
    -0.5f,  0.5f, 0.0f   // 3: top left
};
unsigned int indices[] = {
    0, 1, 3,   // first triangle
    1, 2, 3    // second triangle
};

unsigned int VAO, VBO, EBO;
glGenVertexArrays(1, &VAO);
glGenBuffers(1, &VBO);
glGenBuffers(1, &EBO);

glBindVertexArray(VAO);
glBindBuffer(GL_ARRAY_BUFFER, VBO);
glBufferData(GL_ARRAY_BUFFER, sizeof(vertices), vertices, GL_STATIC_DRAW);

glBindBuffer(GL_ELEMENT_ARRAY_BUFFER, EBO);   // the EBO is recorded into the VAO too
glBufferData(GL_ELEMENT_ARRAY_BUFFER, sizeof(indices), indices, GL_STATIC_DRAW);

glVertexAttribPointer(0, 3, GL_FLOAT, GL_FALSE, 3 * sizeof(float), (void*)0);
glEnableVertexAttribArray(0);

// Drawing: glDrawElements instead of glDrawArrays
glBindVertexArray(VAO);
glDrawElements(GL_TRIANGLES, 6, GL_UNSIGNED_INT, 0); // 6 indices
\end{lstlisting}

\uyari{EBO bağını (\code{GL\_ELEMENT\_ARRAY\_BUFFER}) VAO bağlıyken çözme! VAO,
o an bağlı EBO'yu hatırlar. Eğer VAO'yu çözmeden önce EBO'yu 0'a bağlarsan
çizim başarısız olur. Doğru sıra: önce VAO'yu çöz, sonra istersen EBO'yu.}

\gsub{Uniform'lar: CPU'dan Shader'a Veri}

\emph{Öznitelikler} her köşe için farklıdır; \textbf{uniform}'lar ise tüm çizim
boyunca sabit, CPU'dan gönderilen global değerlerdir. Örnek: rengi zamanla
değiştiren bir üçgen.

\begin{lstlisting}[style=glslstyle]
// FRAGMENT SHADER
#version 330 core
out vec4 FragColor;
uniform vec4 dynamicColor;   // will come from the CPU

void main() {
    FragColor = dynamicColor;
}
\end{lstlisting}

\begin{lstlisting}
// C++ main loop: oscillate the color over time
glUseProgram(shaderProgram);  // use the program FIRST!
float t     = (float)glfwGetTime();
float green = (sin(t) / 2.0f) + 0.5f;   // oscillates in 0..1

int location = glGetUniformLocation(shaderProgram, "dynamicColor");
glUniform4f(location, 0.0f, green, 0.0f, 1.0f);  // update the uniform

glBindVertexArray(VAO);
glDrawArrays(GL_TRIANGLES, 0, 3);
\end{lstlisting}

\uyari{\code{glUniform*} çağrısı, güncellenecek uniform'un ait olduğu programın
\emph{o an kullanımda} olmasını gerektirir. Yani önce \code{glUseProgram}, sonra
\code{glUniform}. \code{glGetUniformLocation} $-1$ dönerse uniform bulunamamış
(veya kullanılmadığı için silinmiş) demektir.}

% =====================================================================
\gsec{Dokular (Textures)}
% =====================================================================

Doku, bir yüzeye "yapıştırdığımız" resimdir. Her köşeye bir \emph{doku koordinatı}
(UV, $[0,1]$ aralığında) veririz; rasterizasyon bunları interpolasyonla yayar ve
fragment shader dokudan örnekleyerek (sample) rengi okur. Resmi yüklemek için
hafif \code{stb\_image.h} kütüphanesi standarttır.

\begin{lstlisting}
#define STB_IMAGE_IMPLEMENTATION
#include "stb_image.h"

unsigned int texture;
glGenTextures(1, &texture);
glBindTexture(GL_TEXTURE_2D, texture);

// Wrapping and filtering settings
glTexParameteri(GL_TEXTURE_2D, GL_TEXTURE_WRAP_S, GL_REPEAT);
glTexParameteri(GL_TEXTURE_2D, GL_TEXTURE_WRAP_T, GL_REPEAT);
glTexParameteri(GL_TEXTURE_2D, GL_TEXTURE_MIN_FILTER, GL_LINEAR_MIPMAP_LINEAR);
glTexParameteri(GL_TEXTURE_2D, GL_TEXTURE_MAG_FILTER, GL_LINEAR);

// Load the image
int width, height, channels;
stbi_set_flip_vertically_on_load(true);   // OpenGL's UV origin is bottom-left!
unsigned char* data = stbi_load("wall.jpg", &width, &height, &channels, 0);
if (data) {
    glTexImage2D(GL_TEXTURE_2D, 0, GL_RGB, width, height, 0,
                 GL_RGB, GL_UNSIGNED_BYTE, data);
    glGenerateMipmap(GL_TEXTURE_2D);
} else {
    std::cerr << "Failed to load texture\n";
}
stbi_image_free(data);   // free the CPU-side data (already copied to the GPU)
\end{lstlisting}

\begin{lstlisting}[style=glslstyle]
// VERTEX SHADER — forward the UV coordinate
#version 330 core
layout (location = 0) in vec3 aPos;
layout (location = 1) in vec2 aUV;
out vec2 uv;
void main() {
    gl_Position = vec4(aPos, 1.0);
    uv = aUV;
}
\end{lstlisting}

\begin{lstlisting}[style=glslstyle]
// FRAGMENT SHADER — sample from the texture
#version 330 core
in  vec2 uv;
out vec4 FragColor;
uniform sampler2D image;   // bound texture unit
void main() {
    FragColor = texture(image, uv);   // read the color at uv
}
\end{lstlisting}

\gsub{İki Dokuyu Karıştırmak (Texture Units)}

Birden fazla doku kullanırken her birini bir \emph{birime} bağlarsın. Klasik
alıştırma: iki dokuyu \code{mix} ile harmanlamak.

\begin{lstlisting}
// Bind two textures to units 0 and 1
glActiveTexture(GL_TEXTURE0);
glBindTexture(GL_TEXTURE_2D, texture0);
glActiveTexture(GL_TEXTURE1);
glBindTexture(GL_TEXTURE_2D, texture1);

// Tell each sampler which unit to read from
glUseProgram(shaderProgram);
glUniform1i(glGetUniformLocation(shaderProgram, "textureA"), 0);
glUniform1i(glGetUniformLocation(shaderProgram, "textureB"), 1);
\end{lstlisting}

\begin{lstlisting}[style=glslstyle]
// FRAGMENT SHADER — blend two textures 80% / 20%
#version 330 core
in  vec2 uv;
out vec4 FragColor;
uniform sampler2D textureA;
uniform sampler2D textureB;
void main() {
    FragColor = mix(texture(textureA, uv), texture(textureB, uv), 0.2);
}
\end{lstlisting}

\uyari{\code{stbi\_set\_flip\_vertically\_on\_load(true)} genelde şarttır: resim
formatları soldan-üstten başlar, OpenGL doku koordinatları ise soldan-alttan.
Çevirmezsen dokular baş aşağı görünür.}

% =====================================================================
\gsec{glm ile Dönüşümleri Uygulamak}
% =====================================================================

Kısım 2'nin matrislerini artık gerçek çizime bağlıyoruz. Bir matrisi shader'a
uniform olarak gönderip vertex shader'da köşeleri dönüştüreceğiz.

\begin{lstlisting}
#include <glm/glm.hpp>
#include <glm/gtc/matrix_transform.hpp>
#include <glm/gtc/type_ptr.hpp>   // for value_ptr

// Inside the main loop:
glm::mat4 transform(1.0f);
transform = glm::translate(transform, glm::vec3(0.5f, -0.5f, 0.0f)); // to bottom-right
transform = glm::rotate(transform, (float)glfwGetTime(),             // spin over time
                        glm::vec3(0.0f, 0.0f, 1.0f));

glUseProgram(shaderProgram);
unsigned int loc = glGetUniformLocation(shaderProgram, "transform");
// send the mat4: 1 matrix, do not transpose (GL_FALSE), data pointer
glUniformMatrix4fv(loc, 1, GL_FALSE, glm::value_ptr(transform));
\end{lstlisting}

\begin{lstlisting}[style=glslstyle]
// VERTEX SHADER
#version 330 core
layout (location = 0) in vec3 aPos;
layout (location = 1) in vec2 aUV;
out vec2 uv;
uniform mat4 transform;
void main() {
    gl_Position = transform * vec4(aPos, 1.0);  // apply the transform
    uv = aUV;
}
\end{lstlisting}

\uyari{\code{glUniformMatrix4fv}'nin üçüncü parametresi \code{transpose}'dur. GLM
matrisleri \emph{sütun-öncelikli} (column-major) saklar ve OpenGL de öyle bekler,
bu yüzden \textbf{GL\_FALSE} verilir. GL\_TRUE verirsen matris transpoze olur ve
dönüşümler bozulur. \code{glm::value\_ptr} ise mat4'ün ham float işaretçisini verir.}

% =====================================================================
\gsec{3B'ye Geçiş: Küp ve Derinlik Testi}
% =====================================================================

Şimdi düz üçgenden gerçek 3B'ye geçiyoruz. MVP zincirini kuracak, derinlik
tamponunu açacak ve dönen bir küp çizeceğiz.

\gsub{Derinlik Testini Açmak}

3B'de arkadaki nesneler öndekiler tarafından örtülmeli. Bunu \textbf{derinlik
testi} (z-buffer) sağlar. Açmazsan üçgenler çizim sırasına göre üst üste biner.

\begin{lstlisting}
// Once, during setup:
glEnable(GL_DEPTH_TEST);

// Every frame, clear the depth buffer ALONG WITH the color buffer:
glClear(GL_COLOR_BUFFER_BIT | GL_DEPTH_BUFFER_BIT);
\end{lstlisting}

\uyari{Derinlik testini açtıysan her karede \code{GL\_DEPTH\_BUFFER\_BIT}'i de
temizlemeyi unutma. Unutursan bir önceki karenin derinlik değerleri kalır ve
nesneler rastgele kaybolur/parçalanır. Bu çok yaygın bir 3B hatasıdır.}

\gsub{MVP ile Dönen Küp}

Küp 36 köşeden oluşur (6 yüz $\times$ 2 üçgen $\times$ 3 köşe). Üç matrisi kurup
shader'a gönderelim.

\begin{lstlisting}
glEnable(GL_DEPTH_TEST);

// Main loop:
glClearColor(0.1f, 0.1f, 0.12f, 1.0f);
glClear(GL_COLOR_BUFFER_BIT | GL_DEPTH_BUFFER_BIT);
glUseProgram(shaderProgram);

// MODEL: spin the cube on its own axis over time
glm::mat4 model = glm::rotate(glm::mat4(1.0f),
                              (float)glfwGetTime() * glm::radians(50.0f),
                              glm::vec3(0.5f, 1.0f, 0.0f));
// VIEW: pull the camera back (+3 on Z, world moves -3)
glm::mat4 view = glm::translate(glm::mat4(1.0f), glm::vec3(0.0f, 0.0f, -3.0f));
// PROJECTION: perspective
glm::mat4 projection = glm::perspective(glm::radians(45.0f),
                                        800.0f / 600.0f, 0.1f, 100.0f);

// Send all three
glUniformMatrix4fv(glGetUniformLocation(shaderProgram, "model"),
                   1, GL_FALSE, glm::value_ptr(model));
glUniformMatrix4fv(glGetUniformLocation(shaderProgram, "view"),
                   1, GL_FALSE, glm::value_ptr(view));
glUniformMatrix4fv(glGetUniformLocation(shaderProgram, "projection"),
                   1, GL_FALSE, glm::value_ptr(projection));

glBindVertexArray(VAO);
glDrawArrays(GL_TRIANGLES, 0, 36);   // 36 vertices = full cube
\end{lstlisting}

\begin{lstlisting}[style=glslstyle]
// VERTEX SHADER — the MVP chain
#version 330 core
layout (location = 0) in vec3 aPos;
layout (location = 1) in vec2 aUV;
out vec2 uv;
uniform mat4 model;
uniform mat4 view;
uniform mat4 projection;
void main() {
    // Right to left: model first, then view, then projection
    gl_Position = projection * view * model * vec4(aPos, 1.0);
    uv = aUV;
}
\end{lstlisting}

\bilgi{İşte Kısım 2'deki $\mathbf{v}_{clip}=\mathbf{P}\cdot\mathbf{V}\cdot\mathbf{M}
\cdot\mathbf{v}$ formülünün canlı hali. Model küpü döndürür, View onu kameranın
önüne taşır, Projection perspektif verir. Bu üç matris 3B grafiğin bel kemiğidir.}

% =====================================================================
\gsec{Kamera Sınıfı (FPS Kamera)}
% =====================================================================

Sabit kamera sıkıcı. Şimdi klavye+fare ile gezilebilen bir FPS kamerası
yazacağız. Euler açılarını (yaw, pitch) kullanıp bir yön vektörü üreteceğiz ve
\code{lookAt} ile view matrisini her karede güncelleyeceğiz.

\begin{lstlisting}
#include <glm/glm.hpp>
#include <glm/gtc/matrix_transform.hpp>

enum class Direction { Forward, Backward, Left, Right };

class Camera {
public:
    glm::vec3 position;
    glm::vec3 front  = glm::vec3(0.0f, 0.0f, -1.0f);
    glm::vec3 up     = glm::vec3(0.0f, 1.0f,  0.0f);
    float yaw        = -90.0f;   // -90: looks toward -Z at start
    float pitch      =  0.0f;
    float speed      =  2.5f;
    float sensitivity = 0.1f;

    explicit Camera(glm::vec3 startPos) : position(startPos) {}

    glm::mat4 viewMatrix() const {
        return glm::lookAt(position, position + front, up);
    }

    // Keyboard: WASD (scaled by deltaTime => FPS-independent)
    void processKeyboard(Direction dir, float deltaTime) {
        float velocity = speed * deltaTime;
        glm::vec3 right = glm::normalize(glm::cross(front, up));
        if (dir == Direction::Forward)  position += front * velocity;
        if (dir == Direction::Backward) position -= front * velocity;
        if (dir == Direction::Left)     position -= right * velocity;
        if (dir == Direction::Right)    position += right * velocity;
    }

    // Mouse: update the look direction
    void processMouse(float xOffset, float yOffset) {
        yaw   += xOffset * sensitivity;
        pitch += yOffset * sensitivity;
        // Clamp the pitch to avoid flipping over
        if (pitch >  89.0f) pitch =  89.0f;
        if (pitch < -89.0f) pitch = -89.0f;
        updateFront();
    }

private:
    void updateFront() {
        // Build a direction vector from yaw/pitch (spherical coordinates)
        glm::vec3 f;
        f.x = cos(glm::radians(yaw)) * cos(glm::radians(pitch));
        f.y = sin(glm::radians(pitch));
        f.z = sin(glm::radians(yaw)) * cos(glm::radians(pitch));
        front = glm::normalize(f);
    }
};
\end{lstlisting}

\begin{lstlisting}
// Usage — in the main loop and callbacks:
Camera camera(glm::vec3(0.0f, 0.0f, 3.0f));

void processInput(GLFWwindow* window, float deltaTime) {
    if (glfwGetKey(window, GLFW_KEY_W) == GLFW_PRESS)
        camera.processKeyboard(Direction::Forward, deltaTime);
    if (glfwGetKey(window, GLFW_KEY_S) == GLFW_PRESS)
        camera.processKeyboard(Direction::Backward, deltaTime);
    if (glfwGetKey(window, GLFW_KEY_A) == GLFW_PRESS)
        camera.processKeyboard(Direction::Left, deltaTime);
    if (glfwGetKey(window, GLFW_KEY_D) == GLFW_PRESS)
        camera.processKeyboard(Direction::Right, deltaTime);
}

// Mouse callback (store lastX/lastY, skip the very first frame):
float lastX = 400.0f, lastY = 300.0f;
bool  firstMouse = true;
void cursorPosCallback(GLFWwindow* win, double xpos, double ypos) {
    if (firstMouse) { lastX = xpos; lastY = ypos; firstMouse = false; }
    float xOffset = xpos - lastX;
    float yOffset = lastY - ypos;   // reversed: screen Y goes down, camera Y goes up
    lastX = xpos; lastY = ypos;
    camera.processMouse(xOffset, yOffset);
}

// In the draw step: get the view matrix from the camera
glm::mat4 view = camera.viewMatrix();
\end{lstlisting}

\ipucu{Bu FPS kamera çapraz çarpımı (sağ yönü bulmak için \code{cross(front, up)}),
küresel koordinatları (yaw/pitch $\to$ yön) ve \code{lookAt}'ı bir arada kullanır.
Kısım 2'deki her araç burada iş başında. Zoom için scroll callback'inde FOV'u
değiştirebilirsin.}

% =====================================================================
\gsec{Temel Aydınlatma: Phong Modeli}
% =====================================================================

Işıksız 3B düz görünür. Phong aydınlatma modeli ışığı üç bileşene ayırır ve
gerçekçi bir hacim hissi verir. Her bileşen Kısım 2'deki vektör işlemlerine
dayanır.

\begin{center}
\begin{tabularx}{\textwidth}{@{}>{\bfseries\color{anaMavi}}l >{\raggedright\arraybackslash}X@{}}
\toprule
\rowcolor{acikMavi}\color{anaMavi}\textbf{Bileşen} & \color{anaMavi}\textbf{Fikir}\\
\midrule
Ambient & Ortam ışığı; her yer minimum aydınlık (sabit çarpan).\\
\rowcolor{acikGri} Diffuse & Yüzeyin ışığa dönüklüğü: $\max(\mathbf{N}\cdot\mathbf{L}, 0)$ (nokta çarpım!).\\
Specular & Parlak yansıma noktası; bakış ve yansıma yönüne bağlı.\\
\bottomrule
\end{tabularx}
\end{center}

\begin{lstlisting}[style=glslstyle]
// FRAGMENT SHADER — Phong lighting
#version 330 core
in  vec3 Normal;    // surface normal (in world space)
in  vec3 FragPos;   // fragment's world position
out vec4 FragColor;

uniform vec3 lightPos;
uniform vec3 lightColor;
uniform vec3 objectColor;
uniform vec3 viewPos;

void main() {
    // 1) Ambient — constant base light
    float ambientStrength = 0.1;
    vec3 ambient = ambientStrength * lightColor;

    // 2) Diffuse — dot product of normal and light direction
    vec3 N = normalize(Normal);
    vec3 L = normalize(lightPos - FragPos);
    float diff = max(dot(N, L), 0.0);
    vec3 diffuse = diff * lightColor;

    // 3) Specular — how much the reflected ray points at the camera
    float specularStrength = 0.5;
    vec3 V = normalize(viewPos - FragPos);
    vec3 R = reflect(-L, N);
    float spec = pow(max(dot(V, R), 0.0), 32.0); // 32 = shininess exponent
    vec3 specular = specularStrength * spec * lightColor;

    // Combine
    vec3 result = (ambient + diffuse + specular) * objectColor;
    FragColor = vec4(result, 1.0);
}
\end{lstlisting}

\begin{lstlisting}[style=glslstyle]
// VERTEX SHADER — forward the normal and world position
#version 330 core
layout (location = 0) in vec3 aPos;
layout (location = 1) in vec3 aNormal;
out vec3 Normal;
out vec3 FragPos;
uniform mat4 model;
uniform mat4 view;
uniform mat4 projection;
void main() {
    FragPos = vec3(model * vec4(aPos, 1.0));       // world position
    // Normal matrix: fixes normals under non-uniform scaling
    Normal  = mat3(transpose(inverse(model))) * aNormal;
    gl_Position = projection * view * model * vec4(aPos, 1.0);
}
\end{lstlisting}

\uyari{Normalleri \code{model} matrisiyle doğrudan çarpma! Nesne \emph{eşit olmayan}
ölçeklenmişse (ör. X'te 2, Y'de 1) normaller yüzeye dik kalmaz. Doğru dönüşüm
\emph{normal matrisidir}: \code{mat3(transpose(inverse(model)))}. Ayrıca ışık ve
bakış yönlerinin \textbf{normalize} olduğundan emin ol.}

\bilgi{Diffuse bileşenindeki $\max(\mathbf{N}\cdot\mathbf{L}, 0)$ tam olarak Kısım
2'de öğrendiğimiz nokta çarpımın $\cos\theta$ anlamıdır: yüzey ışığa ne kadar
dik dönükse o kadar aydınlık. Grafik matematiği ve kodun burada tam olarak
buluştuğunu görüyorsun.}

% #####################################################################
\gpart{5}{Projeler ve İyi Uygulamalar}
% #####################################################################

% =====================================================================
\gsec{Soyutlanmış Shader Sınıfı}
% =====================================================================

Shader derleme kodunu her seferinde yazmak yorucu. Yeniden kullanılabilir bir
sınıf yazalım: dosyadan okur, derler, hata raporlar ve uniform ayarlamayı
kolaylaştırır. Bu, gerçek projelerin standart yapısıdır.

\begin{lstlisting}
#include <glad/glad.h>
#include <glm/glm.hpp>
#include <glm/gtc/type_ptr.hpp>
#include <string>
#include <fstream>
#include <sstream>
#include <iostream>

class Shader {
public:
    unsigned int id;

    Shader(const char* vertexPath, const char* fragmentPath) {
        std::string vertexCode   = readFile(vertexPath);
        std::string fragmentCode = readFile(fragmentPath);
        unsigned int vs = compile(GL_VERTEX_SHADER,   vertexCode.c_str(),   "VERTEX");
        unsigned int fs = compile(GL_FRAGMENT_SHADER, fragmentCode.c_str(), "FRAGMENT");

        id = glCreateProgram();
        glAttachShader(id, vs);
        glAttachShader(id, fs);
        glLinkProgram(id);
        checkLink(id);
        glDeleteShader(vs);
        glDeleteShader(fs);
    }

    void use() const { glUseProgram(id); }

    // Uniform helpers
    void setBool (const std::string& name, bool value)  const { glUniform1i(loc(name), (int)value); }
    void setInt  (const std::string& name, int value)   const { glUniform1i(loc(name), value); }
    void setFloat(const std::string& name, float value) const { glUniform1f(loc(name), value); }
    void setVec3 (const std::string& name, const glm::vec3& v) const {
        glUniform3fv(loc(name), 1, glm::value_ptr(v));
    }
    void setMat4 (const std::string& name, const glm::mat4& m) const {
        glUniformMatrix4fv(loc(name), 1, GL_FALSE, glm::value_ptr(m));
    }

private:
    int loc(const std::string& name) const {
        return glGetUniformLocation(id, name.c_str());
    }

    static std::string readFile(const char* path) {
        std::ifstream file(path);
        if (!file) { std::cerr << "Cannot open file: " << path << "\n"; return ""; }
        std::stringstream ss; ss << file.rdbuf();
        return ss.str();
    }

    static unsigned int compile(GLenum type, const char* src, const char* label) {
        unsigned int shader = glCreateShader(type);
        glShaderSource(shader, 1, &src, nullptr);
        glCompileShader(shader);
        int success; glGetShaderiv(shader, GL_COMPILE_STATUS, &success);
        if (!success) {
            char log[1024]; glGetShaderInfoLog(shader, 1024, nullptr, log);
            std::cerr << label << " shader error:\n" << log << "\n";
        }
        return shader;
    }

    static void checkLink(unsigned int program) {
        int success; glGetProgramiv(program, GL_LINK_STATUS, &success);
        if (!success) {
            char log[1024]; glGetProgramInfoLog(program, 1024, nullptr, log);
            std::cerr << "Program link error:\n" << log << "\n";
        }
    }
};
\end{lstlisting}

\begin{lstlisting}
// See how much this simplifies usage:
Shader shader("shaders/basic.vert", "shaders/basic.frag");

// In the main loop:
shader.use();
shader.setMat4("model",      model);
shader.setMat4("view",       camera.viewMatrix());
shader.setMat4("projection", projection);
shader.setVec3("lightColor", glm::vec3(1.0f));
\end{lstlisting}

\ipucu{Bu Shader sınıfı hemen her OpenGL projesinde bulunur. Aynı mantıkla bir
\code{Mesh} (VAO/VBO/EBO'yu saran), bir \code{Texture} ve \code{Camera} sınıfı
yazarak temiz bir mini motor iskeletine ulaşırsın.}

% =====================================================================
\gsec{Tam Program: Dönen Renkli Küp}
% =====================================================================

Şimdiye kadar öğrendiğimiz her parçayı birleştiren, kopyalayıp derlenebilir tam
bir program. GLFW penceresi + GLAD + derinlik testi + MVP + dönen küp.

\begin{lstlisting}
#include <glad/glad.h>
#include <GLFW/glfw3.h>
#include <glm/glm.hpp>
#include <glm/gtc/matrix_transform.hpp>
#include <glm/gtc/type_ptr.hpp>
#include <iostream>

const char* vertexSrc = R"(
#version 330 core
layout (location = 0) in vec3 aPos;
layout (location = 1) in vec3 aColor;
out vec3 vColor;
uniform mat4 mvp;
void main() {
    gl_Position = mvp * vec4(aPos, 1.0);
    vColor = aColor;
})";

const char* fragmentSrc = R"(
#version 330 core
in vec3 vColor;
out vec4 FragColor;
void main() { FragColor = vec4(vColor, 1.0); })";

void framebufferSizeCallback(GLFWwindow*, int width, int height) {
    glViewport(0, 0, width, height);
}

int main() {
    glfwInit();
    glfwWindowHint(GLFW_CONTEXT_VERSION_MAJOR, 3);
    glfwWindowHint(GLFW_CONTEXT_VERSION_MINOR, 3);
    glfwWindowHint(GLFW_OPENGL_PROFILE, GLFW_OPENGL_CORE_PROFILE);

    GLFWwindow* window = glfwCreateWindow(800, 600, "Spinning Cube", nullptr, nullptr);
    glfwMakeContextCurrent(window);
    glfwSetFramebufferSizeCallback(window, framebufferSizeCallback);
    gladLoadGLLoader((GLADloadproc)glfwGetProcAddress);
    glEnable(GL_DEPTH_TEST);

    // 36 vertices, each face a different color
    float vertices[] = {
        // position           // color
        -0.5f,-0.5f,-0.5f,  1,0,0,   0.5f,-0.5f,-0.5f,  1,0,0,   0.5f, 0.5f,-0.5f,  1,0,0,
         0.5f, 0.5f,-0.5f,  1,0,0,  -0.5f, 0.5f,-0.5f,  1,0,0,  -0.5f,-0.5f,-0.5f,  1,0,0,
        -0.5f,-0.5f, 0.5f,  0,1,0,   0.5f,-0.5f, 0.5f,  0,1,0,   0.5f, 0.5f, 0.5f,  0,1,0,
         0.5f, 0.5f, 0.5f,  0,1,0,  -0.5f, 0.5f, 0.5f,  0,1,0,  -0.5f,-0.5f, 0.5f,  0,1,0,
        -0.5f, 0.5f, 0.5f,  0,0,1,  -0.5f, 0.5f,-0.5f,  0,0,1,  -0.5f,-0.5f,-0.5f,  0,0,1,
        -0.5f,-0.5f,-0.5f,  0,0,1,  -0.5f,-0.5f, 0.5f,  0,0,1,  -0.5f, 0.5f, 0.5f,  0,0,1,
         0.5f, 0.5f, 0.5f,  1,1,0,   0.5f, 0.5f,-0.5f,  1,1,0,   0.5f,-0.5f,-0.5f,  1,1,0,
         0.5f,-0.5f,-0.5f,  1,1,0,   0.5f,-0.5f, 0.5f,  1,1,0,   0.5f, 0.5f, 0.5f,  1,1,0,
        -0.5f,-0.5f,-0.5f,  0,1,1,   0.5f,-0.5f,-0.5f,  0,1,1,   0.5f,-0.5f, 0.5f,  0,1,1,
         0.5f,-0.5f, 0.5f,  0,1,1,  -0.5f,-0.5f, 0.5f,  0,1,1,  -0.5f,-0.5f,-0.5f,  0,1,1,
        -0.5f, 0.5f,-0.5f,  1,0,1,   0.5f, 0.5f,-0.5f,  1,0,1,   0.5f, 0.5f, 0.5f,  1,0,1,
         0.5f, 0.5f, 0.5f,  1,0,1,  -0.5f, 0.5f, 0.5f,  1,0,1,  -0.5f, 0.5f,-0.5f,  1,0,1
    };

    unsigned int VAO, VBO;
    glGenVertexArrays(1, &VAO); glGenBuffers(1, &VBO);
    glBindVertexArray(VAO);
    glBindBuffer(GL_ARRAY_BUFFER, VBO);
    glBufferData(GL_ARRAY_BUFFER, sizeof(vertices), vertices, GL_STATIC_DRAW);
    glVertexAttribPointer(0, 3, GL_FLOAT, GL_FALSE, 6*sizeof(float), (void*)0);
    glEnableVertexAttribArray(0);
    glVertexAttribPointer(1, 3, GL_FLOAT, GL_FALSE, 6*sizeof(float),
                          (void*)(3*sizeof(float)));
    glEnableVertexAttribArray(1);

    // Compile shaders inline (use the Shader class in real projects)
    auto compile = [](GLenum type, const char* src) {
        unsigned int s = glCreateShader(type);
        glShaderSource(s, 1, &src, nullptr); glCompileShader(s); return s;
    };
    unsigned int program = glCreateProgram();
    glAttachShader(program, compile(GL_VERTEX_SHADER,   vertexSrc));
    glAttachShader(program, compile(GL_FRAGMENT_SHADER, fragmentSrc));
    glLinkProgram(program);

    while (!glfwWindowShouldClose(window)) {
        if (glfwGetKey(window, GLFW_KEY_ESCAPE) == GLFW_PRESS)
            glfwSetWindowShouldClose(window, true);

        glClearColor(0.08f, 0.08f, 0.10f, 1.0f);
        glClear(GL_COLOR_BUFFER_BIT | GL_DEPTH_BUFFER_BIT);

        glm::mat4 model = glm::rotate(glm::mat4(1.0f),
            (float)glfwGetTime(), glm::vec3(0.5f, 1.0f, 0.0f));
        glm::mat4 view  = glm::translate(glm::mat4(1.0f), glm::vec3(0, 0, -3));
        glm::mat4 projection = glm::perspective(glm::radians(45.0f),
            800.0f / 600.0f, 0.1f, 100.0f);
        glm::mat4 mvp = projection * view * model;

        glUseProgram(program);
        glUniformMatrix4fv(glGetUniformLocation(program, "mvp"),
                           1, GL_FALSE, glm::value_ptr(mvp));
        glBindVertexArray(VAO);
        glDrawArrays(GL_TRIANGLES, 0, 36);

        glfwSwapBuffers(window);
        glfwPollEvents();
    }

    glDeleteVertexArrays(1, &VAO);
    glDeleteBuffers(1, &VBO);
    glDeleteProgram(program);
    glfwTerminate();
    return 0;
}
\end{lstlisting}

\ipucu{Bu programı derleyip çalıştırdığında, koyu bir arka planda her yüzü farklı
renkte, kendi ekseninde dönen bir küp göreceksin. Buraya kadar öğrendiğin her
kavram --- pencere, bağlam, VBO/VAO, shader, MVP, derinlik testi --- tek bir
çalışan örnekte bir araya geliyor.}

% =====================================================================
\gsec{Kaynak Yönetimi ve Sık Yapılan Hatalar}
% =====================================================================

\gsub{Kaynakları Temizlemek}

OpenGL nesneleri (VAO, VBO, doku, shader) GPU belleğinde durur; program biterken
veya nesne gereksizken \textbf{elle silinmelidir}.

\begin{lstlisting}
glDeleteVertexArrays(1, &VAO);
glDeleteBuffers(1, &VBO);
glDeleteBuffers(1, &EBO);
glDeleteTextures(1, &texture);
glDeleteProgram(shaderProgram);
glfwDestroyWindow(window);
glfwTerminate();
\end{lstlisting}

\gsub{Hata Ayıklama Kontrol Listesi}

\begin{center}
\begin{tabularx}{\textwidth}{@{}>{\bfseries\color{kirmizi}}l >{\raggedright\arraybackslash}X@{}}
\toprule
\rowcolor{kirmizi!12}\color{kirmizi}\textbf{Belirti} & \color{kirmizi}\textbf{Olası neden}\\
\midrule
Siyah/boş ekran & Shader derlenmedi; \code{glUseProgram} yok; VAO bağlı değil; near=0.\\
\rowcolor{acikGri} Segfault (başta) & GLAD yüklenmeden GL çağrısı; bağlam güncel değil.\\
Nesne bozuk/parçalı & Derinlik tamponu temizlenmiyor; derinlik testi kapalı.\\
\rowcolor{acikGri} Her şey ters/aynalı & Matris transpoze (GL\_TRUE) veya yanlış çarpım sırası.\\
Aşırı hızlı/yavaş hareket & Delta time kullanılmıyor.\\
\rowcolor{acikGri} Doku baş aşağı & \code{stbi\_set\_flip\_vertically\_on\_load} çağrılmadı.\\
Aydınlatma yanlış & Normaller/yön vektörleri normalize değil; normal matrisi kullanılmadı.\\
\bottomrule
\end{tabularx}
\end{center}

\bilgi{OpenGL hatalarını yakalamak için \code{glGetError()} kullanabilir ya da
OpenGL 4.3+'ta \code{glDebugMessageCallback} ile otomatik hata mesajları
alabilirsin. Geliştirme sırasında bunu açmak, sessiz hataları görünür kılar ve
saatlerce süren hata avını dakikalara indirir.}

\gsub{Nereden Devam Etmeli?}

\begin{center}
\begin{tabularx}{\textwidth}{@{}>{\bfseries\color{anaMavi}}l >{\raggedright\arraybackslash}X@{}}
\toprule
\rowcolor{acikMavi}\color{anaMavi}\textbf{Konu} & \color{anaMavi}\textbf{Ne öğrenirsin}\\
\midrule
Model yükleme (Assimp) & OBJ/FBX gibi hazır 3B modelleri sahneye almak.\\
\rowcolor{acikGri} Framebuffer'lar & Ekran-dışı çizim, post-processing (bloom, blur).\\
Gölge haritaları & Gerçekçi gölgeler (shadow mapping).\\
\rowcolor{acikGri} Instancing & Binlerce nesneyi tek çağrıda çizmek.\\
PBR & Fizik tabanlı gerçekçi malzemeler.\\
\rowcolor{acikGri} Dear ImGui & GLFW/OpenGL üstüne hızlı arayüz (debug paneli).\\
\bottomrule
\end{tabularx}
\end{center}

\ipucu{Bu rehber seni grafiğin \emph{ne} olduğundan \emph{nasıl} yapıldığına
taşıdı: teori $\to$ matematik $\to$ pencere $\to$ üçgen $\to$ dokulu, aydınlatılmış,
kameralı 3B sahne. Buradan sonrası pratik: küçük sahneler kur, shader'larla oyna,
her hatayı bir öğrenme fırsatı gör. İyi çizimler!}

\vfill
\begin{center}
{\color{ortaGri}\rule{0.5\linewidth}{0.6pt}}\\[6pt]
{\color{ortaGri}\small C++ ile GLFW \& OpenGL Rehberi --- Son}
\end{center}

\end{document}
